% =============================================================
% Bölüm 1: Sayı Sistemleri ve Kodlama
% =============================================================
\chapter{Sayı Sistemleri ve Kodlama}
\label{chap:sayi-sistemleri}

\begin{tcolorbox}[
    colback=blokyesil!15,
    colframe=sinyalyesil!60!black,
    title={\textbf{Öğrenme Hedefleri}},
    fonttitle=\bfseries,
    rounded corners,
    boxrule=0.8pt
]
Bu bölümü tamamladığınızda aşağıdakileri yapabileceksiniz:
\begin{itemize}[nosep, leftmargin=*]
    \item Pozisyonel sayı sistemlerinin genel yapısını ve herhangi bir tabandaki sayıyı
    ondalığa çevirmeyi uygulamak.
    \item Ondalık bir sayıyı ardışık bölme yöntemiyle istenen tabana çevirmek.
    \item İkili, sekizlik ve onaltılık sistemler arasında doğrudan bit gruplamasıyla dönüşüm
    yapmak.
    \item İkili sayılarda toplama, çıkarma, çarpma ve bölme işlemlerini genel $r$-tabanlı
    formülleriyle birlikte elle yapmak.
    \item $(r{-}1)$'in ve $r$'nin tümleyenlerini (9'un/10'un, 1'in/2'nin) hesaplamak ve
    çıkarma işlemini toplamaya çevirmek için kullanmak.
    \item İşaretli tam sayıların (işaret-büyüklük, 1'in ve 2'nin tümleyeni) gösterimlerini
    karşılaştırmak ve 2'nin tümleyeninin neden tercih edildiğini gerekçelendirmek.
    \item 2'nin tümleyeni sistemiyle işaretli toplama/çıkarma yapmak ve taşma (overflow)
    kavramını tanımlamak.
    \item IEEE 754 standardıyla kayan noktalı (ondalıklı) sayıları normalize edip
    işaret/üs/mantissa alanlarına ayrıştırmak, ve $0.1+0.2\neq0.3$ gibi yuvarlama
    tuhaflıklarının kaynağını açıklamak.
    \item BCD, Gray kodu ve ASCII gibi temel ikili kodların kullanım amaçlarını açıklamak.
\end{itemize}
\end{tcolorbox}

\section{Giriş}
\label{sec:sayi-sistemleri-giris}

Bir bilgisayarın belleğinde tuttuğu her şey (bir tam sayı, bir metin karakteri, bir
fotoğrafın bir pikseli, bir ses örneği) sonunda aynı hammaddeye indirgenir: $0$ ve $1$'lerden
oluşan bir dizi. Bu indirgemenin nasıl yapıldığını, yani sayıların ve sembollerin ikili
düzende nasıl \emph{kodlandığını} anlamadan, bu bitleri işleyecek devreleri tasarlamaya
başlayamayız. Bu yüzden yolculuğumuza, tüm sayısal donanım tasarımının en alt katmanı olan
sayı sistemleriyle başlıyoruz: önce herhangi bir tabandaki bir sayıyı nasıl okuyup
dönüştüreceğimizi, ardından negatif sayıları ve metni/sembolleri ikili düzende nasıl
temsil edeceğimizi göreceğiz.

\section{Sayı Tabanları ve Dönüşümler}
\label{sec:sayi-tabanlari}

\subsection{Pozisyonel Sayı Sistemi}
\label{sec:pozisyonel-sistem}

\begin{tanim}{Pozisyonel Sayı Sistemi:}
Bir sayı sisteminde her basamağın değeri, o basamağın \emph{konumuna} bağlıysa buna
\emph{pozisyonel (ağırlıklı) sayı sistemi}\index{pozisyonel sayı sistemi} denir. $r$
\emph{tabanında}\index{taban (radix)} (radix) yazılmış bir sayı, aşağıdaki ağırlıklı toplamla
ifade edilir:
\[
(d_n d_{n-1} \ldots d_1 d_0)_r = \sum_{i=0}^{n} d_i \cdot r^i, \qquad 0 \le d_i < r
\]
Günlük kullandığımız ondalık (decimal) sistem, $r=10$ özel durumudur; örneğin
$137 = 1\cdot10^2+3\cdot10^1+7\cdot10^0$.
\end{tanim}

\begin{bilginot}[Neden Bilgisayarlar İkili Kullanır?]
İkili (binary, $r=2$) sistemin seçimi matematiksel değil, \emph{fiziksel} bir mühendislik
kararıdır. Bir devrede onluk sistem için gereken $10$ farklı gerilim seviyesini güvenilir
biçimde ayırt etmek pahalı ve gürültüye karşı kırılgan olurdu. Seviyeler arası fark küçük
kalır, en ufak bir parazit yanlış okumaya yol açar. Yalnızca \emph{iki} durumu (akım var/yok,
yüksek/düşük gerilim) ayırt etmek çok daha güvenilir ve ucuzdur; seviyeler arasında büyük bir
güvenlik payı (gürültü marjı) bırakılabilir. İkili değişkenler üzerine kurulu Bool cebirini
Bölüm~\ref{chap:bool-cebri}'de işleyeceğiz.
\end{bilginot}

\subsection{İkiliden Ondalığa Dönüşüm}
\label{sec:ikili-ondalik}

Bir tabandan ondalığa dönüşüm, pozisyonel sistem tanımındaki ağırlıklı toplamın doğrudan
uygulanmasıyla yapılır. Zorluğu artan üç örnekle pekiştirelim.

\begin{ornek}[$(101)_2$'nin Ondalık Karşılığı]
\[
(101)_2 = 1\cdot2^2+0\cdot2^1+1\cdot2^0 = 4+0+1 = (5)_{10}
\]
\end{ornek}

\begin{ornek}[$(1011)_2$'nin Ondalık Karşılığı]
\[
(1011)_2 = 1\cdot2^3+0\cdot2^2+1\cdot2^1+1\cdot2^0 = 8+0+2+1 = (11)_{10}
\]
\end{ornek}

\begin{ornek}[$(110101)_2$'nin Ondalık Karşılığı]
\[
(110101)_2 = 1\cdot2^5+1\cdot2^4+0\cdot2^3+1\cdot2^2+0\cdot2^1+1\cdot2^0 = 32+16+0+4+0+1 = (53)_{10}
\]
\end{ornek}

\subsection{Ondalıktan İkiliye Dönüşüm}
\label{sec:ondalik-ikili}

Ondalıktan bir tabana geçiş, ağırlıklı toplamın tersine çevrilmesiyle yapılır: tam sayı kısmı
hedef tabana \emph{ardışık olarak bölünür}; her adımın kalanı, \emph{sondan başa} doğru
okunarak sonuç elde edilir. Şimdi bunu, zorluğu artan üç örnekle, hep ikilik tabanda
pekiştirelim.

\begin{ornek}[$(6)_{10}$'un İkili Karşılığı]
\begin{center}
\begin{tabular}{lll}
\toprule
İşlem & Bölüm & Kalan \\
\midrule
$6 \div 2$ & $3$ & $0$ (en düşük anlamlı bit) \\
$3 \div 2$ & $1$ & $1$ \\
$1 \div 2$ & $0$ & $1$ (en yüksek anlamlı bit) \\
\bottomrule
\end{tabular}
\end{center}
Kalanlar aşağıdan yukarı okunduğunda: $(6)_{10} = (110)_2$.
\end{ornek}

\begin{bilginot}[Neden Kalan, Bölüm Değil?]
$N = d_n r^n+\ldots+d_1 r+d_0$ olsun. $d_0$ dışındaki her terim $r$'nin katı olduğundan,
$N \bmod r = d_0$; yani \textbf{kalan, her zaman tam olarak en düşük anlamlı basamaktır}.
Buna karşılık $\lfloor N/r\rfloor = d_n r^{n-1}+\ldots+d_1$: \textbf{bölüm, aynı sayının son
basamağı silinmiş hâlidir}, kendisi asla bir basamak değildir. İşlemi bölüm $0$'a düşene kadar
tekrarlarız ($0$'ı basamak olarak yazmayız); her adımda yazılan basamak o adımın
\emph{kalanı}dır, bölümü değil. Yukarıdaki örneğin ikinci satırında ($3\div2=1$ kalan $1$)
bölüm ($1$) ile kalan ($1$) aynı değeri taşıyor; bu, satırı yanlış okuyup bölümü basamak
sanmanın fark edilmeyeceği bir durum yaratır. Aşağıdaki iki örnekte bölüm ve kalanın genelde
\emph{birbirinden bağımsız, farklı} sayılar olduğunu göreceğiz.
\end{bilginot}

\begin{ornek}[$(25)_{10}$'un İkili Karşılığı]
\begin{center}
\begin{tabular}{lll}
\toprule
İşlem & Bölüm & Kalan \\
\midrule
$25 \div 2$ & $12$ & $1$ (en düşük anlamlı bit) \\
$12 \div 2$ & $6$ & $0$ \\
$6 \div 2$ & $3$ & $0$ \\
$3 \div 2$ & $1$ & $1$ \\
$1 \div 2$ & $0$ & $1$ (en yüksek anlamlı bit) \\
\bottomrule
\end{tabular}
\end{center}
Kalanlar aşağıdan yukarı okunduğunda: $(25)_{10} = (11001)_2$.
\end{ornek}

\begin{ornek}[$(45)_{10}$'un İkili Karşılığı]
\begin{center}
\begin{tabular}{lll}
\toprule
İşlem & Bölüm & Kalan \\
\midrule
$45 \div 2$ & $22$ & $1$ (en düşük anlamlı bit) \\
$22 \div 2$ & $11$ & $0$ \\
$11 \div 2$ & $5$ & $1$ \\
$5 \div 2$ & $2$ & $1$ \\
$2 \div 2$ & $1$ & $0$ \\
$1 \div 2$ & $0$ & $1$ (en yüksek anlamlı bit) \\
\bottomrule
\end{tabular}
\end{center}
Kalanlar aşağıdan yukarı okunduğunda: $(45)_{10} = (101101)_2$. Doğrulama:
$32+8+4+1=45$. \checkmark
\end{ornek}

\subsection{Sekizlik ve Onaltılık Sistemler: Sembol Sorunu}
\label{sec:sekizlik-onaltilik}

Onluk sistemde tam $10$ rakam sembolümüz vardır: $0,1,\ldots,9$. Bu, $r\le10$ olan her taban
için yeterlidir. Sekizlik sistemde yalnızca $0$-$7$ kullanılır, $8$ ve $9$'a hiç
dokunulmaz. Ama $r=16$ olduğunda her basamak $0$'dan $15$'e kadar \textbf{16 farklı değer}
alabilmelidir; elimizdeki $10$ rakam sembolü $6$ değer için ($10,11,\ldots,15$) yetersiz kalır.

\begin{bilginot}[Neden "10" Değil de Yeni Bir Sembol?]
Akla gelen ilk çözüm, $10$'u iki karakterle ("10") yazmak gibi görünebilir; ama bu,
pozisyonel gösterimin temel kuralını bozar: \textbf{her basamak tek bir karakter olmalıdır},
yoksa gösterim belirsizleşir. $(110)_{16}$ dizisini gördüğünüzde bunun "1, 1, 0" mı yoksa
"1, 10" mı olduğunu ayırt edemezsiniz. Çözüm, onluk rakamların bittiği yerden alfabeyle
devam etmektir: $A=10$, $B=11$, $C=12$, $D=13$, $E=14$, $F=15$. Bu, matematiksel bir
zorunluluk değil, \emph{yalnızca} her basamağı tek karakterde tutmak için benimsenmiş bir
gösterim kuralıdır.
\end{bilginot}

\begin{table}[h]
\centering
\begin{tabular}{ccc}
\toprule
Ondalık & Onaltılık & 4-bit İkili \\
\midrule
0--9 & 0--9 & 0000--1001 \\
10 & A & 1010 \\
11 & B & 1011 \\
12 & C & 1100 \\
13 & D & 1101 \\
14 & E & 1110 \\
15 & F & 1111 \\
\bottomrule
\end{tabular}
\caption{Ondalık--Onaltılık--İkili Karşılık Tablosu}
\label{tab:hex-tablosu}
\end{table}

Bu tablo (Tablo~\ref{tab:hex-tablosu}), donanım tasarımında sürekli kullanılacağı için
ezberlenmeye değer.

\subsection{Sekizlik ve Onaltılıktan Ondalığa Dönüşüm}
\label{sec:sekizlik-onaltilik-ondalik}

Bu dönüşüm, Bölüm~\ref{sec:pozisyonel-sistem}'de kurduğumuz ağırlıklı toplam formülünün doğrudan
uygulanmasıdır. Yeni bir yöntem gerekmez. Onaltılıkta tek fark, bir harf basamağıyla
karşılaştığınızda onu Tablo~\ref{tab:hex-tablosu}'daki sayısal karşılığıyla değiştirmenizdir.

\begin{ornek}[$(247)_8$'in Ondalık Karşılığı]
\[
(247)_8 = 2\cdot8^2+4\cdot8^1+7\cdot8^0 = 128+32+7 = (167)_{10}
\]
\end{ornek}

\begin{ornek}[$(17)_8$'in Ondalık Karşılığı]
\[
(17)_8 = 1\cdot8^1+7\cdot8^0 = 8+7 = (15)_{10}
\]
\end{ornek}

\begin{ornek}[$(3A)_{16}$'nın Ondalık Karşılığı]
Harf basamağı önce sayısal değerine çevrilir: $A=10$.
\[
(3A)_{16} = 3\cdot16^1+10\cdot16^0 = 48+10 = (58)_{10}
\]
\end{ornek}

\begin{ornek}[$(2FC)_{16}$'nın Ondalık Karşılığı]
İki harf basamağı birden çevrilir: $F=15$, $C=12$.
\[
(2FC)_{16} = 2\cdot16^2+15\cdot16^1+12\cdot16^0 = 512+240+12 = (764)_{10}
\]
\end{ornek}

\subsection{Ondalıktan Sekizliğe ve Onaltılığa Dönüşüm}
\label{sec:ondalik-sekizlik-onaltilik}

Bu dönüşüm de yeni bir yöntem gerektirmez: Bölüm~\ref{sec:ondalik-ikili}'de öğrendiğimiz
ardışık bölme yöntemi, taban $2$ yerine $8$ veya $16$ ile aynen uygulanır. Onaltılıkta tek
fark, bir kalan $10$ veya üzerine çıktığında Tablo~\ref{tab:hex-tablosu} ile harfe
çevrilmesidir.

\begin{ornek}[$(100)_{10}$'un Sekizlik Karşılığı]
\begin{center}
\begin{tabular}{lll}
\toprule
İşlem & Bölüm & Kalan \\
\midrule
$100 \div 8$ & $12$ & $4$ (en düşük anlamlı basamak) \\
$12 \div 8$ & $1$ & $4$ \\
$1 \div 8$ & $0$ & $1$ (en yüksek anlamlı basamak) \\
\bottomrule
\end{tabular}
\end{center}
Kalanlar aşağıdan yukarı okunduğunda: $(100)_{10}=(144)_8$. Doğrulama:
$1\cdot8^2+4\cdot8^1+4\cdot8^0=64+32+4=100$. \checkmark
\end{ornek}

\begin{ornek}[$(45)_{10}$'un Sekizlik Karşılığı]
\begin{center}
\begin{tabular}{lll}
\toprule
İşlem & Bölüm & Kalan \\
\midrule
$45 \div 8$ & $5$ & $5$ (en düşük anlamlı basamak) \\
$5 \div 8$ & $0$ & $5$ (en yüksek anlamlı basamak) \\
\bottomrule
\end{tabular}
\end{center}
Kalanlar aşağıdan yukarı okunduğunda: $(45)_{10}=(55)_8$. Doğrulama:
$5\cdot8^1+5\cdot8^0=40+5=45$. \checkmark
\end{ornek}

\begin{ornek}[$(202)_{10}$'un Onaltılık Karşılığı: Bir Tutarlılık Sınaması]
Bu örneği daha önce, Bölüm~\ref{sec:gruplama}'da bit gruplamasıyla da çözüp $(CA)_{16}$
bulacağız. Burada \emph{başka bir yöntemle} aynı sonuca ulaşarak iki yöntemi çapraz
doğrulamış olacağız.
\begin{center}
\begin{tabular}{lll}
\toprule
İşlem & Bölüm & Kalan \\
\midrule
$202 \div 16$ & $12$ & $10=\text{A}$ (en düşük anlamlı basamak) \\
$12 \div 16$ & $0$ & $12=\text{C}$ (en yüksek anlamlı basamak) \\
\bottomrule
\end{tabular}
\end{center}
Kalanlar aşağıdan yukarı okunduğunda: $(202)_{10}=(CA)_{16}$.
\end{ornek}

\begin{ornek}[$(500)_{10}$'un Onaltılık Karşılığı]
\begin{center}
\begin{tabular}{lll}
\toprule
İşlem & Bölüm & Kalan \\
\midrule
$500 \div 16$ & $31$ & $4$ (en düşük anlamlı basamak) \\
$31 \div 16$ & $1$ & $15=\text{F}$ \\
$1 \div 16$ & $0$ & $1$ (en yüksek anlamlı basamak) \\
\bottomrule
\end{tabular}
\end{center}
Kalanlar aşağıdan yukarı okunduğunda: $(500)_{10}=(1F4)_{16}$. Doğrulama:
$1\cdot16^2+15\cdot16^1+4\cdot16^0=256+240+4=500$. \checkmark
\end{ornek}

\subsection{İkili, Sekizlik ve Onaltılık Arasında Gruplama}
\label{sec:gruplama}

\begin{tanim}{İkili-Sekizlik-Onaltılık Arası Gruplama:}
İkili sayılar bilgisayar için idealdir ama insan gözüyle takip edilmesi zahmetlidir. Sekizlik
(octal, $r=8$)\index{sekizlik sistem} ve onaltılık (hexadecimal, $r=16$)\index{onaltılık sistem}
sistemler tam bu yüzden kullanılır: $8=2^3$ ve $16=2^4$ olduğu için, ikili ile bu tabanlar
arasındaki dönüşüm bir bölme/çarpma \emph{gerektirmez}. Yalnızca bitlerin gruplanmasıyla
yapılır: her \textbf{3 bit} bir sekizlik basamağa, her \textbf{4 bit} bir onaltılık basamağa
karşılık gelir.
\end{tanim}

\begin{ornek}[İkiliden Onaltılığa Gruplama]
$(11001010)_2$ sayısını onaltılığa çevirelim; sağdan başlayarak 4'erli gruplarız:
\[
\underbrace{1100}_{\text{C}} \; \underbrace{1010}_{\text{A}} \;\Rightarrow\; (CA)_{16}
\]
\end{ornek}

\begin{ornek}[Onaltılıdan İkiliye Açma: Ters Yön]
Gruplama tersine de çevrilebilir: her onaltılık basamak, kendi 4-bit ikili karşılığıyla
değiştirilir. $(CA)_{16}$'yı ikiliye açalım:
\[
\underbrace{\text{C}}_{1100} \; \underbrace{\text{A}}_{1010} \;\Rightarrow\; (11001010)_2
\]
Bu, bir önceki örneğin tam tersidir; aynı sayı çifti arasında gidip gelebildiğimizi gösterir.
\end{ornek}

\begin{ornek}[İkiliden Sekizliğe Gruplama]
$(101101110)_2$ sayısını sekizliğe çevirelim; bu kez bit sayısı ($9$) tam olarak $3$'e
bölündüğü için baştan sıfır eklemeye gerek yok, doğrudan sağdan 3'erli gruplanır:
\[
\underbrace{101}_{5} \; \underbrace{101}_{5} \; \underbrace{110}_{6} \;\Rightarrow\; (556)_8
\]
Doğrulama, her iki tarafı da onluk tabana çevirip karşılaştırarak yapılır:
$(556)_8 = 5\cdot8^2+5\cdot8^1+6\cdot8^0=320+40+6=(366)_{10}$, ve
$(101101110)_2=1\cdot256+1\cdot64+1\cdot32+1\cdot8+1\cdot4+1\cdot2=(366)_{10}$. İki taraf da
aynı onluk değere ($366$) eşit olduğundan sonuç doğrulanmış olur. \checkmark
\end{ornek}

\begin{alistirma}[\kolay]
$(11001010)_2$ sayısını, yukarıdaki örnekteki gibi ama bu kez \textbf{sekizlik} sisteme,
sağdan 3'erli gruplayarak çeviriniz. (İpucu: 8 bit 3'e tam bölünmez; en sola gerektiği kadar
$0$ ekleyerek grupları tamamlayınız.)
\end{alistirma}

\noindent\textbf{Çözüm:} 8 biti 3'e tam bölmek için sola bir $0$ eklenir, 9 bit elde edilir;
sağdan 3'erli gruplanır:
\[
\underbrace{011}_{3} \; \underbrace{001}_{1} \; \underbrace{010}_{2} \;\Rightarrow\; (312)_8
\]
Doğrulama: $(312)_8 = 3\cdot8^2+1\cdot8^1+2\cdot8^0 = 192+8+2=202$, ve
$(11001010)_2=(202)_{10}$ --- iki sonuç örtüşüyor.

\subsection{Tabanı Bilinmeyen İşlemlerde Taban Bulma}
\label{sec:bilinmeyen-taban}

Şimdiye kadar hep taban $r$'yi bilip bir sayıyı çeviriyorduk. Bu alt bölümde işi tersine
çevireceğiz: pozisyonel sistem tanımındaki ağırlıklı toplam formülünü ($\sum d_i r^i$),
tabanın kendisini \textbf{bilinmeyen} kabul ederek bir denklem olarak kuracak ve $r$ için
çözeceğiz. Tek kısıtlama akılda tutulmalı: bir ifadede kullanılan her basamak, tabandan
\emph{küçük} olmak zorundadır ($d_i < r$); bulunan çözüm bu kısıtla sınanmalıdır.

\begin{ornek}[$(23)_x$ Sayısının Ondalık Karşılığı $19$ İse $x$ Nedir?]
Ağırlıklı toplam formülü, üsler açık yazılarak kurulur; $x^0=1$ olduğu için birler basamağı
$x$ ile çarpılmaz, sadece kendisi kalır:
\[
(23)_x = 2\cdot x^1+3\cdot x^0 = 2x+3
\]
Bu ifade $19$'a eşitlenip $x$ için çözülür:
\[
2x+3=19 \;\Rightarrow\; 2x=16 \;\Rightarrow\; x=8
\]
Kısıt kontrolü: kullanılan basamaklar $\{2,3\}$, $x=8$ bunların ikisinden de büyük. \checkmark
\end{ornek}

\begin{ornek}[$(34)_x+(21)_x=(105)_x$ İse $x$ Nedir?]
Üç tarafı da önce üsler açık yazılarak kuralım:
\[
(3\cdot x^1+4\cdot x^0)+(2\cdot x^1+1\cdot x^0) = 1\cdot x^2+0\cdot x^1+5\cdot x^0
\]
Sadeleştirilince:
\[
(3x+4)+(2x+1) = x^2+5 \;\Rightarrow\; 5x+5 = x^2+5 \;\Rightarrow\; x^2-5x=0
\;\Rightarrow\; x(x-5)=0
\]
İki kök: $x=0$ ve $x=5$. $x=0$ geçerli bir taban olamaz (ayrıca kısıt $x>4$'ü de sağlamaz);
$x=5$ kısıtı sağlar ($5>4$). \checkmark Doğrulama: $(34)_5{=}19$, $(21)_5{=}11$, toplam
$30=(105)_5$.
\end{ornek}

\begin{ornek}[$(32)_x+(24)_x=(100)_x$ İse $x$ Nedir?]
Bu kez eşitliğin sağ tarafı (right-hand side, RHS) üç basamaklı olduğu için denklem
\emph{ikinci dereceden} çıkar. Üsler açık:
\[
(3\cdot x^1+2\cdot x^0)+(2\cdot x^1+4\cdot x^0) = 1\cdot x^2+0\cdot x^1+0\cdot x^0
\]
Sadeleştirilince:
\[
(3x+2)+(2x+4) = x^2 \;\Rightarrow\; x^2-5x-6=0 \;\Rightarrow\; (x-6)(x+1)=0
\]
İki kök: $x=6$ ve $x=-1$. Negatif taban fiziksel olarak anlamsız olduğundan $x=-1$ reddedilir;
$x=6$ kısıtı da sağlar ($6>4$). \checkmark Doğrulama: $(32)_6{=}20$, $(24)_6{=}16$, toplam
$36=(100)_6$.
\end{ornek}

\subsection{Kesirli Sayı Dönüşümleri}
\label{sec:kesirli-donusum}

Pozisyonel gösterim yalnızca tam sayılarla sınırlı değildir; noktadan sonraki basamaklar
\emph{negatif} kuvvetler taşır:
\[
(0.d_{-1}d_{-2}\ldots d_{-m})_r = \sum_{i=1}^{m} d_{-i}\cdot r^{-i}
\]
Tabandan ondalığa dönüşüm yine doğrudan bu toplamla yapılır; örneğin
$(0.101)_2 = 1\cdot2^{-1}+0\cdot2^{-2}+1\cdot2^{-3} = 0.5+0+0.125 = (0.625)_{10}$. Ondalıktan
bir tabana geçiş ise, tam sayı kısmından \textbf{farklı ve ters işleyen} bir yöntem gerektirir.

\begin{ornek}[Ondalıktan Bir Tabana Dönüşüm: Ardışık Çarpma Yöntemi]
Kesirli kısım hedef tabanla ardışık olarak çarpılır; her adımda taşan tam sayı basamağı
sonucun bir hanesi olur ve kalan kesirle işleme devam edilir. $(0.6875)_{10}$'u ikiliye
çevirelim:
\begin{center}
\begin{tabular}{lll}
\toprule
İşlem & Sonuç & Taşan Basamak \\
\midrule
$0.6875\times2$ & $1.375$ & $1$ (en yüksek anlamlı bit) \\
$0.375\times2$ & $0.75$ & $0$ \\
$0.75\times2$ & $1.5$ & $1$ \\
$0.5\times2$ & $1.0$ & $1$ (en düşük anlamlı bit) \\
\bottomrule
\end{tabular}
\end{center}
Kesir $0$'a düştüğü için işlem sonlanır. Taşan basamaklar \textbf{yukarıdan aşağı} okunur:
$(0.6875)_{10} = (0.1011)_2$.
\end{ornek}

\begin{bilginot}[Tuzak: Yön Simetrik Değildir]
Tam sayı dönüşümünde (ardışık bölme) sonuç \emph{aşağıdan yukarı} okunurken, kesir
dönüşümünde (ardışık çarpma) sonuç \emph{yukarıdan aşağı} okunur. Bu asimetri sık yapılan
bir hatanın kaynağıdır. Bir sayının tam ve kesirli kısımlarını ayrı ayrı çevirip doğru
yönde birleştirmeyi unutmayın.
\end{bilginot}

\begin{bilginot}[Her Kesir Sonlanmaz]
Her ondalık kesir, hedef tabanda sonlu basamakla ifade edilemez. Bunu somut görmek için
$(0.1)_{10}$'u ardışık çarpma yöntemiyle ikiliye çevirmeyi deneyelim:
\begin{center}
\begin{tabular}{lll}
\toprule
İşlem & Sonuç & Taşan Basamak \\
\midrule
$0.1\times2$ & $0.2$ & $0$ \\
$0.2\times2$ & $0.4$ & $0$ \\
$0.4\times2$ & $0.8$ & $0$ \\
$0.8\times2$ & $1.6$ & $1$ \\
$0.6\times2$ & $1.2$ & $1$ \\
$0.2\times2$ & $0.4$ & $0$ \\
\bottomrule
\end{tabular}
\end{center}
Altıncı satırda kesir tekrar $0.2$'ye döner; bu değeri ikinci satırda zaten görmüştük.
Yani işlem hiçbir zaman $0$'a düşmeyen bir \textbf{döngüye} girmiştir: kesir sonsuza kadar
$0.2\to0.4\to0.8\to0.6\to0.2\to\ldots$ sırasını tekrar eder ve her turda aynı $0,0,1,1$
basamak dörtlüsünü üretir. Sonuç, sonsuz tekrar eden bir ikili dizidir:
$(0.1)_{10}=(0.0001100110011\ldots)_2$. Bu olgu, kayan noktalı sayı gösteriminde
karşılaşılan $0.1+0.2\neq0.3$ gibi tanıdık tuhaflıkların temel nedenidir: sonlu bit
genişliğinde bu sonsuz dizi bir yerde kesilmek zorunda kalır, ve kesilen kısım küçük bir
yuvarlama hatası olarak kalır. Konuya kayan nokta gösterimini işlerken geri döneceğiz.
\end{bilginot}

\begin{alistirma}[\kolay]
$(0.375)_{10}$'u ardışık çarpma yöntemiyle ikiliye çeviriniz.
\end{alistirma}

\noindent\textbf{Çözüm:}
\begin{center}
\begin{tabular}{lll}
\toprule
İşlem & Sonuç & Taşan Basamak \\
\midrule
$0.375\times2$ & $0.75$ & $0$ \\
$0.75\times2$ & $1.5$ & $1$ \\
$0.5\times2$ & $1.0$ & $1$ \\
\bottomrule
\end{tabular}
\end{center}
Yukarıdan aşağı okunduğunda: $(0.375)_{10}=(0.011)_2$. Doğrulama:
$0\cdot2^{-1}+1\cdot2^{-2}+1\cdot2^{-3}=0.25+0.125=0.375$. \checkmark

\begin{ornek}[Sekizlik Kesir Dönüşümü]
Ardışık çarpma yöntemi herhangi bir tabana genelleşir; $8$'lik tabana dönüşümde kesirli
kısım $2$ yerine $8$ ile çarpılır. $(0.513)_{10}$'u
sekizliğe çevirelim:
\begin{center}
\begin{tabular}{lll}
\toprule
İşlem & Sonuç & Taşan Basamak \\
\midrule
$0.513\times8$ & $4.104$ & $4$ \\
$0.104\times8$ & $0.832$ & $0$ \\
$0.832\times8$ & $6.656$ & $6$ \\
$0.656\times8$ & $5.248$ & $5$ \\
$0.248\times8$ & $1.984$ & $1$ \\
$0.984\times8$ & $7.872$ & $7$ \\
\bottomrule
\end{tabular}
\end{center}
Kesir hiç $0$'a düşmüyor (beklendiği gibi --- her kesir sonlanmaz); $6$ basamak hassasiyetle
kesersek: $(0.513)_{10} \approx (0.406517)_8$.
\end{ornek}

\begin{ornek}[Karma Sayı Dönüşümü: Tam ve Kesirli Kısım Birlikte]
Hem tam sayı hem kesir içeren bir sayıyı çevirirken, \textbf{iki kısım ayrı ayrı çevrilip
sonra birleştirilir}; tam kısım için ardışık bölme, kesirli kısım için ardışık çarpma.
$(41.6875)_{10}$'u ikiliye çevirelim.

Tam kısım ($41$), ardışık bölme ile:
\begin{center}
\begin{tabular}{lll}
\toprule
İşlem & Bölüm & Kalan \\
\midrule
$41 \div 2$ & $20$ & $1$ (en düşük anlamlı bit) \\
$20 \div 2$ & $10$ & $0$ \\
$10 \div 2$ & $5$ & $0$ \\
$5 \div 2$ & $2$ & $1$ \\
$2 \div 2$ & $1$ & $0$ \\
$1 \div 2$ & $0$ & $1$ (en yüksek anlamlı bit) \\
\bottomrule
\end{tabular}
\end{center}
Kalanlar aşağıdan yukarı: $(41)_{10}=(101001)_2$.

Kesirli kısım ($0.6875$), ardışık çarpma ile bulunur (bunu daha önce hesaplamıştık):
$(0.6875)_{10}=(0.1011)_2$.

İki sonuç birleştirilir:
\[
(41.6875)_{10} = (101001.1011)_2
\]
\end{ornek}

\section{İkili Aritmetik: Toplama, Çıkarma, Çarpma ve Bölme}
\label{sec:ikili-aritmetik}

Şimdiye kadar sayıları temsil etmeyi ve tabanlar arasında çevirmeyi öğrendik. Bu bölümde bu
sayılar üzerinde temel aritmetik işlemleri (toplama, çıkarma, çarpma, bölme) nasıl
yapacağımızı ele alıyoruz. $r$ tabanındaki aritmetik işlemler, ilkokulda öğrendiğimiz ondalık
kurallarla \textbf{birebir aynı mantıkla} işler; tek fark, her basamağın $0$'dan $r{-}1$'e kadar
sınırlı olması ve elde/borç eşiğinin $10$ değil $r$ olmasıdır. Bir sonraki bölümde (Tümleyenler)
çıkarmayı toplamaya çeviren bir "trik" öğreneceğiz. Ama önce, optimize edeceğimiz ham
işlemi elle nasıl yaptığımızı bilmemiz gerekiyor. Her işlem için önce genel $r$-tabanlı
matematiksel formülü kuracağız, sonra $r=2$ özel durumunu (ikili) örneklerle pekiştireceğiz;
bu formüller, ileride mantık kapılarıyla gerçek toplayıcı/çarpıcı/bölücü \emph{devreleri}
tasarlarken birebir kullanılacak.

\subsection{İkili Toplama}
\label{sec:ikili-toplama}

\begin{tanim}{Genel Toplama Formülü:}
$n$ basamaklı, $r$ tabanındaki $A=(a_{n-1}\ldots a_1a_0)_r$ ile $B=(b_{n-1}\ldots b_1b_0)_r$
toplanırken, her basamak $i$ için (sağdan sola, $i{=}0$'dan başlayarak):
\[
s_i = (a_i+b_i+c_i) \bmod r, \qquad c_{i+1} = \left\lfloor \frac{a_i+b_i+c_i}{r} \right\rfloor
\]
Burada $c_i$, bir önceki basamaktan gelen elde ($c_0{=}0$). $a_i,b_i,c_i \le r-1$ olduğundan
toplamları en fazla $2(r{-}1)+1=2r{-}1<2r$'dir; bu yüzden $c_{i+1}$ her zaman $0$ veya $1$'dir.
Hiçbir zaman $2$ veya daha büyük bir elde oluşmaz.
\end{tanim}

İkilide ($r{=}2$) bu formül $s_i=(a_i+b_i+c_i)\bmod2$, $c_{i+1}=\lfloor(a_i+b_i+c_i)/2\rfloor$'ye
indirgenir; bu, ileride kuracağımız bir "tam toplayıcı" (full adder) mantık kapısı
devresinin giriş-çıkış ilişkisinin ta kendisidir. Ondalıkta iki basamağın toplamı $10$'u geçince
elde doğar; ikilide eşik $r{=}2$'dir: $1+1=(10)_2$ --- yani bu basamağa $0$ yazılır, bir üst
basamağa $1$ elde gider.

\begin{ornek}[$(101101)_2+(100111)_2$]
Sağdan sola, sütun sütun:
\begin{center}
\begin{tabular}{cccc}
\toprule
Basamak & İşlem & Yaz & Elde \\
\midrule
0 & $1+1$ & $0$ & $1$ \\
1 & $0+1+1$ & $0$ & $1$ \\
2 & $1+1+1$ & $1$ & $1$ \\
3 & $1+0+1$ & $0$ & $1$ \\
4 & $0+0+1$ & $1$ & $0$ \\
5 & $1+1+0$ & $0$ & $1$ \\
\bottomrule
\end{tabular}
\end{center}
Son elde ($1$) yeni bir basamak olarak yazılır: $(1010100)_2$. Doğrulama:
$(101101)_2=45$, $(100111)_2=39$, $45+39=84=(1010100)_2$. \checkmark Bu bit-bit prosedür,
donanımda birebir bir \emph{tam toplayıcı} (full adder) devresine dönüşür. Devresini
ilerideki Kombinasyonel Mantık bölümünde kuracağız.
\end{ornek}

\begin{alistirma}[\kolay]
$(1101)_2+(1011)_2$ işlemini sütun sütun (elde takip ederek) yapınız.
\end{alistirma}

\noindent\textbf{Çözüm:}
\begin{center}
\begin{tabular}{cccc}
\toprule
Basamak & İşlem & Yaz & Elde \\
\midrule
0 & $1+1$ & $0$ & $1$ \\
1 & $0+1+1$ & $0$ & $1$ \\
2 & $1+0+1$ & $0$ & $1$ \\
3 & $1+1+1$ & $1$ & $1$ \\
\bottomrule
\end{tabular}
\end{center}
Son elde ($1$) yeni basamak olur: $(11000)_2$. Doğrulama: $13+11=24=(11000)_2$. \checkmark

\subsection{İkili Çıkarma}
\label{sec:ikili-cikarma}

\begin{tanim}{Genel Çıkarma Formülü:}
$n$ basamaklı $A-B$ işlemi için, her basamak $i$'de:
\[
d_i = \begin{cases}
a_i-k_i-b_i & \text{eğer } a_i-k_i \ge b_i \\
(a_i-k_i-b_i)+r & \text{eğer } a_i-k_i < b_i
\end{cases}
\qquad
k_{i+1} = \begin{cases}
0 & \text{eğer } a_i-k_i \ge b_i \\
1 & \text{eğer } a_i-k_i < b_i
\end{cases}
\]
Burada $k_i$, bir önceki basamaktan gelen borç ($k_0{=}0$). İkinci durumda ($a_i-k_i<b_i$),
mevcut basamağa taban $r$ \textbf{eklenerek} (yani "komşudan $r$ birim ödünç alınarak") fark
hesaplanır. Ondalıkta bu eklenen değer $10$, ikilide $2$'dir.
\end{tanim}

Ondalıkta borç alırken basamağa $10$ eklenir; ikilide borç alırken basamağa sadece $2$
eklenir. Genel kural: \textbf{borcun değeri, her zaman kullanılan tabana eşittir}.

\begin{ornek}[$(101101)_2-(100111)_2$]
\begin{center}
\begin{tabular}{cccc}
\toprule
Basamak & İşlem & Yaz & Borç? \\
\midrule
0 & $1-1$ & $0$ & hayır \\
1 & $0-1 \to (0+2)-1$ & $1$ & evet, üstten alındı \\
2 & $(1-1)-1 \to (0+2)-1$ & $1$ & evet \\
3 & $(1-1)-0$ & $0$ & hayır \\
4 & $0-0$ & $0$ & hayır \\
5 & $1-1$ & $0$ & hayır \\
\bottomrule
\end{tabular}
\end{center}
Sonuç: $(000110)_2$. Doğrulama: $45-39=6=(000110)_2$. \checkmark
\end{ornek}

\begin{alistirma}[\kolay]
$(1101)_2-(1011)_2$ işlemini sütun sütun (borç takip ederek) yapınız.
\end{alistirma}

\noindent\textbf{Çözüm:}
\begin{center}
\begin{tabular}{cccc}
\toprule
Basamak & İşlem & Yaz & Borç? \\
\midrule
0 & $1-1$ & $0$ & hayır \\
1 & $0-1 \to (0+2)-1$ & $1$ & evet \\
2 & $(1-1)-0$ & $0$ & hayır \\
3 & $(1-0)-1$ & $0$ & hayır \\
\bottomrule
\end{tabular}
\end{center}
Sonuç: $(0010)_2$. Doğrulama: $13-11=2=(0010)_2$. \checkmark

\subsection{İkili Çarpma}
\label{sec:ikili-carpma}

\begin{tanim}{Genel Çarpma Formülü:}
$A$ (çarpılan) ile $n$ basamaklı $B=(b_{n-1}\ldots b_1b_0)_r$ (çarpan) çarpımı, çarpanın her
basamağının ürettiği bir \emph{kısmi çarpım}ın, doğru yere kaydırılıp toplanmasıyla bulunur:
\[
A\times B = \sum_{i=0}^{n-1} b_i \cdot A \cdot r^i
\]
\end{tanim}

İkilide ($r{=}2$) çarpan basamakları $b_i$ yalnızca $0$ veya $1$ olabileceğinden, her terim
($b_i\cdot A\cdot 2^i$) ya tamamen \textbf{sıfır} ya da $A$'nın $i$ basamak
\textbf{sola kaydırılmış bir kopyası}dır. Çarpma işlemi hiçbir zaman gerçek bir çarpım
gerektirmez, sadece kaydırma ve toplama. Bu yüzden ikili çarpmaya \emph{kaydır-ve-topla}
(shift-and-add) da denir.

\begin{ornek}[$(1011)_2\times(101)_2$]
Çarpanın basamakları (sağdan sola) $b_0{=}1,\ b_1{=}0,\ b_2{=}1$. Her biri için kısmi çarpımı,
\textbf{zaten kaydırılmış tam genişliğiyle} yazalım (böylece toplama sıradan basamak-hizalı
toplamaya döner):
\begin{center}
\begin{tabular}{lll}
\toprule
Terim & Değer & Açıklama \\
\midrule
$b_0\cdot A\cdot 2^0$ & $001011$ & $A$, kaydırmasız \\
$b_1\cdot A\cdot 2^1$ & $000000$ & $b_1{=}0$ olduğu için tamamen sıfır \\
$b_2\cdot A\cdot 2^2$ & $101100$ & $A$, $2$ basamak sola kaydırılmış \\
\bottomrule
\end{tabular}
\end{center}
Üç terim de artık aynı genişlikte ($6$ bit); Bölüm~\ref{sec:ikili-toplama}'daki sıradan
ikili toplama ile eklenir:
\[
001011 + 000000 + 101100 = 110111
\]
Doğrulama: $(1011)_2=11$, $(101)_2=5$, $11\times5=55=(110111)_2$. \checkmark
\end{ornek}

\subsection{İkili Bölme}
\label{sec:ikili-bolme}

Bölme, diğer üç işlemden farklıdır: basamak başına kapalı bir formülle değil, adım adım
işleyen bir \textbf{algoritma}yla tanımlanır (uzun bölme). Önce ne aradığımızı netleştirelim:

\begin{tanim}{Bölmenin Tanımı:}
$A$ (bölünen) ve $B\neq0$ (bölen) verildiğinde, bölme işlemi
\[
A = Q\cdot B + R, \qquad 0 \le R < B
\]
eşitliğini sağlayan $Q$ (bölüm) ve $R$ (kalan) değerlerini bulmaktır.
\end{tanim}

\begin{tanim}{Genel Bölme Algoritması (Kaydır-ve-Çıkar):}
$A$'nın basamakları en anlamlıdan (soldan) başlayarak tek tek işlenir. Bir çalışma kaydı
$R$ ($0$'dan başlar) her adımda önce tabana göre kaydırılıp bir sonraki basamak eklenir,
sonra bölenle karşılaştırılır:
\[
R \leftarrow r\cdot R + a_i \qquad\text{ardından}\qquad
q_i = \max\{\,q \in \{0,1,\ldots,r{-}1\} : q\cdot B \le R\,\}, \qquad R \leftarrow R - q_i\cdot B
\]
Yani her adımda, $B$'nin kaydırılmış $R$'yi aşmayan en büyük katı bulunur; bu katın kaç
kez alındığı ($q_i$) o adımın bölüm basamağıdır, ve o kat $R$'den çıkarılır. Onluk gibi
$r>2$ tabanlarda $q_i$ genelde $0$ ile $r{-}1$ arasında herhangi bir değer olabileceği için
bu en büyük katı bulmak bir miktar deneme gerektirir (aşağıda somut bir örnekle göreceğiz).
İkilide ($r{=}2$) ise $q_i$ yalnızca $0$ veya $1$ olabileceğinden bu arayış tek bir soruya
indirgenir: \textbf{"şu anki $R$, bölenden büyük veya eşit mi?"} Cevap evetse bölüm biti
$1$ ve $B$ çıkarılır, hayırsa bölüm biti $0$. Bu basitlik (kaydır, karşılaştır, gerekirse
çıkar), donanımda \emph{ardışık bölücü} devrelerinin temelini oluşturur.
\end{tanim}

Önce genel algoritmayı, alışık olduğumuz onluk tabanda (uzun bölme) somutlaştıralım.

\begin{ornek}[$(138)_{10} \div (6)_{10}$: Genel Algoritma Onluk Tabanda]
$A=138$, $B=6$. Basamaklar soldan sağa ($1,3,8$) tek tek işlenir; her adımda $q_i$, $q_i\cdot6$
$R$'yi aşmayacak şekilde \emph{deneme yoluyla} bulunur:
\begin{center}
\begin{tabular}{ccccc}
\toprule
Adım & Basamak & $R\leftarrow10R+a_i$ & En büyük $q\cdot6\le R$ & Bölüm biti, yeni $R$ \\
\midrule
1 & $1$ & $1$ & $0\cdot6\le1$ & $q{=}0$, $R{=}1$ \\
2 & $3$ & $13$ & $2\cdot6\le13$ & $q{=}2$, $R{=}1$ \\
3 & $8$ & $18$ & $3\cdot6\le18$ & $q{=}3$, $R{=}0$ \\
\bottomrule
\end{tabular}
\end{center}
Bölüm basamakları sırayla $0,2,3$: $Q=(023)_{10}=23$. Son $R=0$, kalan yok. Doğrulama:
$23\times6=138$. \checkmark
\end{ornek}

Şimdi aynı algoritmayı, karşılaştırmanın tek bir soruya indirgendiği $r=2$ özel durumunda
uygulayalım:

\begin{ornek}[$(1011)_2 \div (11)_2$]
$A=1011$ ($11$), $B=11$ ($3$). Bitler soldan sağa işlenir:
\begin{center}
\begin{tabular}{ccccc}
\toprule
Adım & Getirilen bit & $R\leftarrow 2R+a_i$ & $R\ge B$ mi? & Bölüm biti, yeni $R$ \\
\midrule
1 & $1$ & $1$ & $1<3$, hayır & $q{=}0$, $R{=}1$ \\
2 & $0$ & $2$ & $2<3$, hayır & $q{=}0$, $R{=}2$ \\
3 & $1$ & $5$ & $5\ge3$, evet & $q{=}1$, $R{=}5-3=2$ \\
4 & $1$ & $5$ & $5\ge3$, evet & $q{=}1$, $R{=}5-3=2$ \\
\bottomrule
\end{tabular}
\end{center}
Bölüm bitleri sırayla $0,0,1,1$: $Q=(0011)_2=3$. Son $R=2=(10)_2$, kalan. Doğrulama:
$11 = 3\times3+2$. \checkmark
\end{ornek}

\begin{alistirma}[\orta]
$(1110)_2 \div (11)_2$ işlemini yukarıdaki tabloyla aynı yöntemle yapınız.
\end{alistirma}

\noindent\textbf{Çözüm:} $A=1110$ ($14$), $B=11$ ($3$):
\begin{center}
\begin{tabular}{ccccc}
\toprule
Adım & Getirilen bit & $R\leftarrow 2R+a_i$ & $R\ge B$ mi? & Bölüm biti, yeni $R$ \\
\midrule
1 & $1$ & $1$ & $1<3$, hayır & $q{=}0$, $R{=}1$ \\
2 & $1$ & $3$ & $3\ge3$, evet & $q{=}1$, $R{=}3-3=0$ \\
3 & $1$ & $1$ & $1<3$, hayır & $q{=}0$, $R{=}1$ \\
4 & $0$ & $2$ & $2<3$, hayır & $q{=}0$, $R{=}2$ \\
\bottomrule
\end{tabular}
\end{center}
$Q=(0100)_2=4$, kalan $R=(10)_2=2$. Doğrulama: $14=4\times3+2$. \checkmark

\section{Tümleyenler}
\label{sec:tumleyenler}

Şimdiye kadar sayıları bir tabandan diğerine çevirdik. Bu bölümde farklı bir amaca hizmet eden
bir araç kuracağız: \emph{tümleyenler}.

\begin{bilginot}[Neden Çıkarmayı Toplamaya Çeviriyoruz?]
Bir bilgisayarın işlemcisi, aslında milyonlarca küçük anahtardan (transistörden) örülmüş bir
ağdır. Bu transistörler, ikili girişlerden ikili çıkış üreten \emph{mantık kapıları}na (AND,
OR, NOT vb.; Bölüm~\ref{sec:mantik-kapilari}'de göreceğimiz kapılara) gruplanır, kapılar da
bir araya gelerek \emph{toplayıcı} (adder) gibi aritmetik devreleri oluşturur. Sorun şu:
toplama işlemini yapan bir devre ile çıkarma işlemini yapan bir devre, \textbf{birbirinden
farklı} kapı düzenlemeleri gerektirir; yani ikisini de donanımda gerçeklemek isterseniz iki
ayrı devre kurmanız, dolayısıyla iki katı transistör harcamanız gerekir. Ama eğer $M-N$
işlemini $M+(\text{tümleyen}(N))$ olarak yeniden yazabilirsek, çıkarma işlemi de aynı toplayıcı
devre üzerinden yürür. Ayrı bir "çıkarıcı" devre tasarlamaya hiç gerek kalmaz. Tümleyenlerin
donanım tasarımındaki asıl değeri budur: transistör/kapı sayısını neredeyse yarıya indiren
köklü bir basitleştirme. Bu devrenin somut kapı şemasını, toplayıcı devreleri işlediğimizde
(ilerideki Toplayıcı Devreler bölümünde) göreceğiz.
\end{bilginot}

\subsection{Genel Tümleyen Tanımı}
\label{sec:tumleyen-tanim}

\begin{tanim}{Tümleyen Türleri:}
$n$ basamaklı, $r$ tabanında bir $N$ sayısı için iki tür tümleyen tanımlanır:
\begin{itemize}[nosep, leftmargin=*]
    \item \textbf{$(r{-}1)$'in tümleyeni}\index{r-1'in tümleyeni} (diminished radix
    complement --- ikilide "1'in tümleyeni", ondalıkta "9'un tümleyeni"):
    \[
    (r{-}1)\text{'in tümleyeni} = (r^n-1) - N
    \]
    \item \textbf{$r$'nin tümleyeni}\index{r'nin tümleyeni} (radix complement --- ikilide
    "2'nin tümleyeni", ondalıkta "10'un tümleyeni"):
    \[
    r\text{'nin tümleyeni} = r^n - N = \big[(r{-}1)\text{'in tümleyeni}\big] + 1
    \]
\end{itemize}
İkinci tanımın ilk tanıma bağlı olması pratik açıdan önemlidir: $(r{-}1)$'in tümleyenini
hesaplamak çok kolaydır (aşağıda göreceğiz), $r$'nin tümleyeni de ona sadece $1$ eklenerek
bulunur.
\end{tanim}

\subsection{Ondalık Örneklerle Tümleyen Hesabı}
\label{sec:tumleyen-ondalik}

Kavramı önce aşina olduğumuz ondalık sistemde (r=10) somutlaştıralım. $(r{-}1)$'in
tümleyenini, yani $9$'un tümleyenini, hesaplamanın pratik yolu: \textbf{her basamağı $9$'dan
çıkarmaktır}; çünkü $(10^n-1)$ sayısı $n$ tane $9$'dan oluşur ($999\ldots9$), ve
basamak basamak çıkarma hiçbir borç gerektirmez.

\begin{ornek}[$546700$'ün 9'un ve 10'un Tümleyeni]
$6$ basamaklı $546700$ sayısının 9'un tümleyeni, her basamağı $9$'dan çıkararak:
\[
9{-}5,\ 9{-}4,\ 9{-}6,\ 9{-}7,\ 9{-}0,\ 9{-}0 \;\Rightarrow\; 453299
\]
Doğrulama: $(10^6-1)-546700 = 999999-546700=453299$. \checkmark

10'un tümleyeni, 9'un tümleyenine $1$ eklenerek bulunur:
\[
453299+1=453300
\]
Doğrulama: $10^6-546700=453300$. \checkmark
\end{ornek}

\begin{ornek}[$012398$'in 9'un Tümleyeni]
Baştaki sıfır da bir basamaktır, aynı kurala tabidir:
\[
9{-}0,\ 9{-}1,\ 9{-}2,\ 9{-}3,\ 9{-}9,\ 9{-}8 \;\Rightarrow\; 987601
\]
\end{ornek}

\subsection{Kısayol Yöntemi ve İspatı}
\label{sec:tumleyen-kisayol}

\begin{bilginot}[Pratik Kısayol: $r$'nin Tümleyenini Doğrudan Almak]
$r$'nin tümleyenini iki adımda (önce $(r{-}1)$'in tümleyenini alıp sonra $1$ eklemek yerine)
\textbf{tek adımda} bulmanın bir kısayolu vardır: sağdan başlayarak \textbf{sondaki sıfırları
ve ilk sıfır-olmayan basamağı olduğu gibi bırakın}, geri kalan (soldaki) basamakları
$(r{-}1)$'den çıkarın.
\end{bilginot}

\begin{ornek}[Kısayolla $246700$'ün 10'un Tümleyeni]
Sondaki iki $0$ aynı kalır; ilk sıfır olmayan basamak ($7$) $10$'dan çıkarılır
($10{-}7=3$); geri kalan basamaklar ($2,4,6$, soldan sağa) $9$'dan çıkarılır ($7,5,3$):
\[
246700 \;\Rightarrow\; 753300
\]
Doğrulama: $10^6-246700=753300$. \checkmark
\end{ornek}

\begin{bilginot}[Kısayol Neden İşliyor? Borç Yayılımı İspatı]
$r$'nin tümleyeni, $r^n$ (yani $n$ basamaklı $1\underbrace{00\ldots0}_{n}$ sayısı) ile $N$
arasındaki farktır: $r^n-N$. Bu çıkarmayı elle, sağdan sola borç alarak yapalım. $N$'nin
sondan itibaren $k$ tane sıfırı olsun. İlk $k$ basamakta $0-0$ işlemi yapılır. Borca hiç
gerek yoktur, sonuç $0$ kalır. Basamak $k$'de (sondan itibaren ilk sıfır-olmayan basamak,
değeri $d$) işlem $0-d$'dir; bu borç gerektirir, sonuçta $(r-d)$ elde edilir ve bir borç sola,
basamak $k{+}1$'e, yayılır. Basamak $k{+}1$'de $r^n$'nin kendi basamağı da $0$
olduğundan, gelen borç yüzünden işlem aslında $(0-1)-d_{k+1} = (r-1)-d_{k+1}$'dir (yine bir
borç daha sola yayılarak); bu, en soldaki basamağa kadar devam eder. Yani ilk
sıfır-olmayan basamağın solundaki \textbf{her} basamak, gelen borç nedeniyle otomatik olarak
$(r{-}1)$'den çıkarılmış olur. Kısayol, bu borç yayılımının doğrudan bir sonucudur.
\end{bilginot}

\subsection{İkilide Tümleyenler}
\label{sec:tumleyen-ikili}

Aynı tanımlar $r=2$ için doğrudan uygulanır. $(r{-}1)=1$ olduğundan, \textbf{1'in tümleyeni,
her biti ters çevirmektir} (NOT); çünkü $1-d_i$ işlemi tam olarak biti tersine çevirir
($1-0=1$, $1-1=0$). \textbf{2'nin tümleyeni} ise 1'in tümleyenine $1$ eklenerek bulunur.

\begin{ornek}[$(1011000)_2$'nin 1'in ve 2'nin Tümleyeni]
Her bit ters çevrilir:
\[
(1011000)_2 \;\xrightarrow{\text{NOT}}\; (0100111)_2 \quad \text{(1'in tümleyeni)}
\]
Bu sonuca $1$ eklenir:
\[
0100111+1=0101000 \quad \text{(2'nin tümleyeni)}
\]
\end{ornek}

\begin{ornek}[Kısayolla $(1011000)_2$'nin 2'nin Tümleyeni]
Sağdan başlayarak sondaki sıfırları ($000$) ve ilk $1$'i olduğu gibi bırakırız ($...1000$),
geri kalan sol taraf ($101$) ters çevrilir ($010$):
\[
(1011000)_2 \;\Rightarrow\; (0101000)_2
\]
Yukarıdaki uzun yöntemle bulduğumuz sonuçla aynı. \checkmark Bu kısayol, elle hesap yaparken
(ör. bir sonraki bölümdeki çıkarma örneklerinde) ciddi zaman kazandırır.
\end{ornek}

\subsection{Sekizlik ve Onaltılıkta $(r{-}1)$'in Tümleyeni}
\label{sec:tumleyen-sekizlik-onaltilik}

Aynı mantık her tabana genelleşir: sekizlikte $(r{-}1)=7$ olduğundan, $(r{-}1)$'in tümleyeni
\textbf{her basamağı $7$'den çıkararak} bulunur; onaltılıkta $(r{-}1)=15={\rm F}$ olduğundan,
\textbf{her basamak $F$'den çıkarılır} (harf basamaklarını önce Tablo~\ref{tab:hex-tablosu}
ile sayıya çevirmeyi, sonucu da gerekirse tekrar harfe çevirmeyi unutmayın).

\begin{ornek}[$(5643)_8$'ün 7'nin Tümleyeni]
Her basamak $7$'den çıkarılır:
\[
7{-}5,\ 7{-}6,\ 7{-}4,\ 7{-}3 \;\Rightarrow\; 2134
\]
Yani $(5643)_8$'ün $7$'nin tümleyeni $(2134)_8$'dir. Doğrulama: $(7777)_8-(5643)_8$ ondalıkta
$4095-2979=1116$, ve $1116$'nın sekizlik karşılığı $(2134)_8$'tür. \checkmark
\end{ornek}

\begin{ornek}[$(3A9)_{16}$'un $F$'nin Tümleyeni]
Her basamak (harfler önce sayıya çevrilerek) $F(=15)$'ten çıkarılır: $15{-}3=12={\rm C}$,
$15{-}10={\rm A}\to5$, $15{-}9=6$:
\[
(3A9)_{16} \;\Rightarrow\; ({\rm C}56)_{16}
\]
Doğrulama: $(FFF)_{16}-(3A9)_{16}$ ondalıkta $4095-937=3158$, ve $3158$'in onaltılık karşılığı
$({\rm C}56)_{16}$'dır. \checkmark
\end{ornek}

\subsection{Kesirli (Noktalı) Sayılarda Tümleyen Alma}
\label{sec:tumleyen-noktali}

Buraya kadarki tüm tanımlar, sayının bir radix noktası (ondalık/ikili nokta) \emph{içermediğini}
varsaydı. Eğer $N$'nin bir noktası varsa, tümleyeni almadan önce nokta \textbf{geçici olarak
kaldırılır}, sayı sanki tam sayıymış gibi tümleyeni alınır, ve nokta \textbf{aynı göreli
konumuna geri konur}.

\begin{ornek}[$(546.700)$'ün 9'un Tümleyeni]
Nokta geçici olarak kaldırılır: $546700$. Bunun 9'un tümleyeni (daha önce hesaplamıştık):
$453299$. Nokta, sondan üçüncü basamağın önündeki aynı göreli konumuna geri konur:
\[
(546.700) \;\Rightarrow\; (453.299)
\]
\end{ornek}

\begin{ornek}[$(101.11)_2$'nin 2'nin Tümleyeni]
Nokta kaldırılır: $10111$ ($5$ basamak). 2'nin tümleyeni: 1'in tümleyeni $01000$, $+1=01001$.
Nokta, sondan iki basamağın önündeki aynı göreli konuma geri konur:
\[
(101.11)_2 \;\Rightarrow\; (010.01)_2
\]
Yöntem tam sayılarla tamamen aynıdır. Tek fark, işlem bitince noktanın eski yerine
konmasıdır.
\end{ornek}

\begin{tanim}{Tümleyenin Tümleyeni:}
$r$'nin tümleyeni tanımı gereği $r^n-N$ olduğundan, bu sonucun bir kez daha $r$'nin
tümleyenini almak
\[
r^n-(r^n-N) = N
\]
verir; yani \textbf{tümleyenin tümleyeni, sayının kendisini geri verir}. Bu, ilerideki
çıkarma prosedüründe "sonuç negatifse, bulduğunuz değerin bir kez daha tümleyenini alıp
başına eksi koyun" adımının neden doğru sonucu verdiğini açıklar: negatif durumda elde
ettiğimiz ara sonuç, aslında gerçek büyüklüğün tümleyenidir; tümleyenini tekrar almak bizi
gerçek büyüklüğe geri götürür.
\end{tanim}

\subsection{$r$'nin Tümleyeniyle Çıkarma İşlemi}
\label{sec:cikarma-r-tumleyeni}

Tümleyenlerin asıl amacına gelelim: $M-N$ çıkarma işlemini toplamaya çevirmek. $n$ basamaklı
$M$ ve $N$ için kural şöyledir:

\begin{tanim}{$r$'nin Tümleyeniyle Çıkarma Kuralı:}
\begin{enumerate}[nosep, leftmargin=*]
    \item $M + [r\text{'nin tümleyeni}(N)]$ toplamını hesaplayın.
    \item Sonuç $n$ basamaktan \textbf{taşarsa} (bir "uç elde" oluşursa), taşan basamağı
    \textbf{atın}. Kalan $n$ basamak doğrudan $M-N$'dir ve sonuç \textbf{pozitiftir}
    ($M \ge N$ demektir).
    \item Sonuç $n$ basamaktan \textbf{taşmazsa}, elde ettiğiniz $n$ basamaklı sonuç aslında
    $N-M$'nin $r$'nin tümleyenidir ($M<N$ demektir). Gerçek büyüklüğü bulmak için bu
    sonucun bir kez daha $r$'nin tümleyenini alın ve başına eksi işareti koyun.
\end{enumerate}
\end{tanim}

\begin{ornek}[Kuralın Doğrudan Uygulanışı: $(48)_{10}-(25)_{10}$]
$n=2$ basamak, $M=48$, $N=25$.

\noindent\textbf{Adım 1:} $M+[10\text{'un tümleyeni}(N)]$ toplamını hesaplayın.
\[
10\text{'un tümleyeni}(25) = 100-25=75
\]
\[
M+\text{tümleyen}(N) = 48+75=123
\]

\noindent\textbf{Adım 2:} Sonuç $2$ basamaktan taşıyor mu diye bakın. $123$'ün en soldaki
$1$'i bir "uç elde"dir, yani sonuç $2$ basamaktan taşmıştır.

\noindent\textbf{Adım 3:} Taştığı için, taşan basamak atılır. Kalan $23$ doğrudan
$M-N$'dir:
\[
48-25=23 \quad \checkmark
\]
\end{ornek}

\begin{bilginot}[Önce Basamak Sayılarını Eşitleyin]
Bu kural, $M$ ve $N$'nin \textbf{aynı sayıda basamağa ($n$)} sahip olmasını varsayar. Tümleyen
tanımı da ($r^n-N$) zaten belirli bir $n$'ye göre yapılır. Sayıların basamak sayıları farklıysa,
kısa olanın \textbf{başına sıfır eklenerek} uzun olanla eşitlenir. Aşağıdaki ilk örnekte
$N=3250$'nin $4$ basamağı var ama $M=72532$'nin $5$; bu yüzden $N$'yi $03250$ olarak yazıp
$n=5$ üzerinden işlem yapıyoruz.
\end{bilginot}

\begin{ornek}[Ondalıkta Çıkarma: Taşma Var, Pozitif Sonuç]
$72532-03250$ işlemini ($n=5$ basamak) tümleyenle yapalım. $N=03250$'nin 10'un tümleyeni:
$100000-03250=96750$.
\[
M+\text{tümleyen}(N) = 72532+96750 = 169282
\]
Bu, $5$ basamaklı sistemde bir "uç elde" (taşan $1$) üretir. Taşan basamak atılır:
\[
72532-3250 = 69282
\]
Doğrulama: $72532-3250=69282$. \checkmark
\end{ornek}

\begin{ornek}[Ondalıkta Çıkarma: Taşma Yok, Negatif Sonuç]
Aynı sayılarla ters yönde, $03250-72532$ işlemini yapalım. $N=72532$'nin 10'un tümleyeni:
$100000-72532=27468$.
\[
M+\text{tümleyen}(N) = 03250+27468 = 30718
\]
Bu $5$ basamağa sığıyor, uç elde \textbf{yok}. Demek ki gerçek sonuç negatif: $30718$'in
10'un tümleyenini alıp başına eksi koyuyoruz:
\[
100000-30718=69282 \;\Rightarrow\; \text{sonuç} = -69282
\]
Doğrulama: $3250-72532=-69282$. \checkmark
\end{ornek}

\begin{ornek}[İkilide Çıkarma: 2'nin Tümleyeniyle, Taşma Var]
$8$ bit üzerinde $(01010100)_2{=}84$'ten $(01000011)_2{=}67$'yi çıkaralım. $N$'nin 2'nin
tümleyeni (kısayolla): $10111101$.
\[
01010100+10111101 = \underbrace{1}_{\text{uç elde}}00010001
\]
Uç elde atılır: sonuç $(00010001)_2=17$. Doğrulama: $84-67=17$. \checkmark
\end{ornek}

\begin{ornek}[İkilide Çıkarma: 2'nin Tümleyeniyle, Taşma Yok]
Ters yönde, $67-84$ işlemini yapalım. $84=(01010100)_2$'nin 2'nin tümleyeni: $10101100$.
\[
01000011+10101100 = 11101111
\]
$8$ bite sığıyor, uç elde yok. Sonuç negatif. $(11101111)_2$'nin 2'nin tümleyeni:
$00010001=17$. Sonuç: $-17$. Doğrulama: $67-84=-17$. \checkmark
\end{ornek}

\subsection{$(r{-}1)$'in Tümleyeniyle Çıkarma: Uçtan-Dolaşan Elde}
\label{sec:cikarma-r-1-tumleyeni}

$(r{-}1)$'in tümleyeni (9'un veya 1'in tümleyeni) ile çıkarma yapmak isterseniz, kural bir
noktada farklılaşır: uç elde oluştuğunda bu sefer \textbf{atılmaz}, bunun yerine sonucun en
düşük anlamlı basamağına \textbf{eklenir}. Buna \emph{uçtan-dolaşan elde} (end-around
carry) denir. Bunun nedeni, $(r{-}1)$'in tümleyeninin $r$'nin tümleyeninden tam $1$ eksik
olması: kaybolan bu farkı telafi etmek için taşan biti geri "dolaştırıyoruz".

\begin{ornek}[Ondalıkta 9'un Tümleyeniyle Çıkarma]
Aynı $72532-3250$ işlemi, bu kez 9'un tümleyeniyle. $N=03250$'nin 9'un tümleyeni: $96749$.
\[
M+\text{tümleyen}(N) = 72532+96749 = 169281
\]
Uç elde ($1$) atılmaz, en düşük basamağa eklenir:
\[
69281+1=69282
\]
Aynı sonuç ($r$'nin tümleyeni yöntemiyle bulduğumuzla özdeş). \checkmark
\end{ornek}

\begin{ornek}[İkilide 1'in Tümleyeniyle Çıkarma]
$84-67$ işlemini bu kez 1'in tümleyeniyle yapalım. $67=(01000011)_2$'nin 1'in tümleyeni:
$10111100$.
\[
01010100+10111100 = \underbrace{1}_{\text{uç elde}}00010000
\]
Uç elde atılmaz, en düşük basamağa eklenir: $00010000+1=00010001=17$. Doğrulama:
$84-67=17$. \checkmark
\end{ornek}

\begin{alistirma}[\orta]
$67-84$ işlemini $8$ bit üzerinde \textbf{1'in tümleyeni} yöntemiyle (uçtan-dolaşan elde
kuralına dikkat ederek) çözünüz.
\end{alistirma}

\noindent\textbf{Çözüm:} $84=(01010100)_2$'nin 1'in tümleyeni: $10101011$.
\[
01000011+10101011 = 11101110
\]
$8$ bite sığıyor, uç elde \textbf{yok}. Bu durumda uçtan-dolaşan elde kuralı devreye
girmez, sonuç doğrudan negatiftir. $(11101110)_2$'nin 1'in tümleyeni: $00010001=17$. Sonuç:
$-17$. Doğrulama: $67-84=-17$. \checkmark

\section{İşaretli İkili Sayılar}
\label{sec:isaretli-sayilar}

Şimdiye kadar gördüğümüz her sayı \emph{işaretsizdi}; yalnızca sıfır ve pozitif değerleri
temsil ediyordu. Ama gerçek dünyada negatif sayılara da ihtiyacımız var. Sıradan
aritmetikte bunu bir "$-$" sembolüyle çözeriz; donanımda ise elimizde yalnızca bitler var,
sembol diye bir şey yok. Bu bölümde, bir önceki bölümde kurduğumuz tümleyen aracını
kullanarak, negatif sayıları saf bitlerle nasıl temsil edeceğimizi göreceğiz.

\subsection{İşaret Biti ve Yorumlama}
\label{sec:isaret-biti}

\begin{tanim}{İşaret Biti:}
İşaretli bir ikili sayıda, en soldaki bit \emph{işaret bitidir}: $0$ pozitif, $1$ negatif
anlamına gelir. Kalan bitler büyüklüğü (ya da tümleyen sistemindeyse, aşağıda göreceğimiz
gibi farklı bir şeyi) temsil eder.
\end{tanim}

\begin{bilginot}[Aynı Bitler, Farklı Anlam]
Bir bit dizisinin işaretli mi işaretsiz mi olduğu, dizinin kendisinde yazılı değildir;
\textbf{yorumlayan tarafından bilinmesi gerekir}. Örneğin $01001$: işaretsiz okunursa $9$,
işaretli okunursa $+9$'dur (soldaki bit $0$ olduğu için ikisi de aynı büyüklüğü verir). Ama
$11001$: işaretsiz okunursa $25$, işaretli okunursa $-9$'dur. En soldaki $1$ işaret biti
olarak yorumlandığında, kalan $1001$ ($9$) büyüklüğünün negatifini gösterir.
\end{bilginot}

\subsection{İşaret-Büyüklük Gösterimi}
\label{sec:isaret-buyukluk}

En sezgisel yöntem, sıradan aritmetikle birebir aynıdır: işaret biti ayrı, büyüklük ayrı
saklanır.

\begin{ornek}[$-5$'in 8 Bitte İşaret-Büyüklük Gösterimi]
$+5=00000101$. İşaret-büyüklükte $-5$, sadece işaret bitini $1$ yaparak elde edilir:
\[
-5 = 10000101
\]
\end{ornek}

\subsection{İşaretli Tümleyen Sistemi: 1'in ve 2'nin Tümleyeni}
\label{sec:isaretli-tumleyen}

İşaret-büyüklük sezgisel olsa da, donanımda kullanılan asıl yöntem farklıdır: negatif
sayı, pozitif sayının \textbf{tümleyeni} alınarak elde edilir; işaret biti dahil,
\textbf{sayının tamamının} tümleyeni.

\begin{ornek}[$-9$'un 8 Bitte Üç Gösterimi]
$+9=00001001$ ($8$ bite tamamlamak için işaret bitinden sonra gerektiği kadar $0$ eklenir).
\begin{center}
\begin{tabular}{ll}
\toprule
Gösterim & $-9$ \\
\midrule
İşaret-büyüklük (sadece işaret biti değişir) & $10001001$ \\
İşaretli 1'in tümleyeni (tüm bitler ters çevrilir) & $11110110$ \\
İşaretli 2'nin tümleyeni (1'in tümleyenine $+1$) & $11110111$ \\
\bottomrule
\end{tabular}
\end{center}
Tablonun ikinci ve üçüncü satırı arasında dikkat edilmesi gereken bir nokta var: 1'in
tümleyenine ($11110110$) eklenen $+1$, işaret bitini değiştirmemiştir; her iki satırda da
en soldaki bit $1$'dir. Bunun nedeni, $+1$ eklemenin sağdan başlayan bir elde zinciri
üretmesi ve bu zincirin (burada $0\to1$ dönüşümüyle) daha ilk basamakta sönmesidir; işaret
bitine hiç ulaşmaz. Sayının negatif olduğunu gösteren işaret zaten bir önceki adımda, tüm
bitlerin (işaret biti dahil) ters çevrilmesiyle oluşmuştu; $+1$ eklemek yalnızca büyüklüğü
küçük bir miktar düzeltir, işareti bozmaz.
\end{ornek}

\begin{bilginot}[İpucu: Ağırlıklı Toplamla Doğrudan Değer Okuma]
$2$'nin tümleyeni gösterimindeki bir sayının değerini bulmak için her zaman "ters çevir,
$1$ ekle, başına eksi koy" yapmak gerekmez. Doğrudan bir kısayol vardır: $n$-bitlik sayı
$b_{n-1}\ldots b_1b_0$ için, \textbf{en soldaki (işaret) bit negatif ağırlıklı} sayılırsa,
değer doğrudan ağırlıklı toplamla okunur:
\[
V = -b_{n-1}\cdot2^{n-1} + \sum_{i=0}^{n-2}b_i\cdot2^i
\]
\end{bilginot}

\begin{ornek}[Kısayolla $(11011)_2$'nin Değerinin Bulunması]
İşaret biti ($b_4=1$) negatif ağırlıklı sayılır:
\[
-1\times2^4+1\times2^3+0\times2^2+1\times2^1+1\times2^0 = -16+8+0+2+1=-5
\]
Yani $(11011)_2=-5$ (2'nin tümleyeni gösteriminde).
\end{ornek}

Bu kısayolun neden çalıştığını görelim. Bölüm~\ref{sec:tumleyenler}'de $r$'nin tümleyenini
$r^n-N$ olarak tanımlamıştık. $n$-bit 2'nin tümleyeninde negatif bir $V$ değeri, bit
kalıbı $X=2^n+V$ olan işaretsiz bir sayı olarak saklanır (yani $V=X-2^n$). Bit kalıbının
işaretsiz değeri $X=\sum_{i=0}^{n-1}b_i2^i$'dir; işaret biti $b_{n-1}=1$ olduğundan
$X=2^{n-1}+\sum_{i=0}^{n-2}b_i2^i$'dir. Buradan:
\[
V=X-2^n=2^{n-1}+\sum_{i=0}^{n-2}b_i2^i-2^n
\]
\[
=-2^{n-1}+\sum_{i=0}^{n-2}b_i2^i
\]
Bu, tam olarak yukarıdaki kısayol formülüdür.

\begin{ornek}[Kısayolun Doğrulanması: $-9$'un $2$'nin Tümleyeni Gösterimi]
Az önce $-9$'un $8$-bit $2$'nin tümleyeni gösterimini $11110111$ olarak bulmuştuk. Aynı
kısayolu bu bit kalıbına uygulayalım:
\[
-1\times2^7+1\times2^6+1\times2^5+1\times2^4+0\times2^3+1\times2^2+1\times2^1+1\times2^0
\]
\[
=-128+64+32+16+0+4+2+1=-9
\]
Doğrulandı.
\end{ornek}

\begin{alistirma}[\kolay]
Kısayolu kullanarak $(10101)_2$'nin ($5$-bit $2$'nin tümleyeni) değerini bulunuz.
\end{alistirma}

\noindent\textbf{Çözüm:} Bitler soldan sağa $b_4b_3b_2b_1b_0=10101$, yani
$b_4{=}1,\,b_3{=}0,\,b_2{=}1,\,b_1{=}0,\,b_0{=}1$. En soldaki (işaret) bit $b_4$ negatif
ağırlıklı sayılarak kısayol formülü kurulur:
\[
V = -b_4\cdot2^4+b_3\cdot2^3+b_2\cdot2^2+b_1\cdot2^1+b_0\cdot2^0
\]
Bit değerleri yerine konur:
\[
V = -1\times2^4+0\times2^3+1\times2^2+0\times2^1+1\times2^0
\]
Her terim ayrı ayrı hesaplanır:
\[
V = -16+0+4+0+1
\]
Toplanır:
\[
V = -11
\]

\subsection{Neden 2'nin Tümleyeni Tercih Edilir}
\label{sec:neden-2-tumleyeni}

\begin{table}[h]
\centering
\begin{tabular}{cccc}
\toprule
Ondalık & 2'nin Tümleyeni & 1'in Tümleyeni & İşaret-Büyüklük \\
\midrule
$+7$ & 0111 & 0111 & 0111 \\
$+6$ & 0110 & 0110 & 0110 \\
$+5$ & 0101 & 0101 & 0101 \\
$+4$ & 0100 & 0100 & 0100 \\
$+3$ & 0011 & 0011 & 0011 \\
$+2$ & 0010 & 0010 & 0010 \\
$+1$ & 0001 & 0001 & 0001 \\
$+0$ & 0000 & 0000 & 0000 \\
$-0$ & --- & 1111 & 1000 \\
$-1$ & 1111 & 1110 & 1001 \\
$-2$ & 1110 & 1101 & 1010 \\
$-3$ & 1101 & 1100 & 1011 \\
$-4$ & 1100 & 1011 & 1100 \\
$-5$ & 1011 & 1010 & 1101 \\
$-6$ & 1010 & 1001 & 1110 \\
$-7$ & 1001 & 1000 & 1111 \\
$-8$ & 1000 & --- & --- \\
\bottomrule
\end{tabular}
\caption{4-Bit İşaretli Sayılar: Üç Gösterimin Karşılaştırması}
\label{tab:isaretli-4bit}
\end{table}

Tablo~\ref{tab:isaretli-4bit}'ten iki gözlem çıkar:

\begin{bilginot}[Sebep 1: Tek Sıfır]
İşaret-büyüklük ve 1'in tümleyeninde \textbf{hem $+0$ hem $-0$} var. $16$ kombinasyonun
biri israf ediliyor, ayrıca "sıfır mı değil mi" karşılaştırması iki farklı bit dizisini
kontrol etmeyi gerektiriyor. 2'nin tümleyeninde sıfırın \textbf{tek} bir gösterimi var; bu
yüzden bir negatif değer daha ($-8$) sığdırılabiliyor.
\end{bilginot}

\begin{bilginot}[Sebep 2 (Asıl Önemli Olan): Aritmetik Kolaylığı]
İşaret-büyüklükte iki sayıyı toplamak için önce işaretlerini \textbf{karşılaştırmanız}
gerekir: işaretler aynıysa büyüklükleri toplayıp ortak işareti verirsiniz; işaretler
farklıysa küçük büyüklüğü büyükten \textbf{çıkarıp} büyük olanın işaretini verirsiniz. Bu,
ayrı bir karşılaştırma \emph{ve} çıkarma devresi gerektirir. 2'nin tümleyeninde ise kural tek
cümledir: \textbf{işaret biti dahil her şeyi normal topla, işaret basamağından taşan biti
at.} Karşılaştırma yok, özel durum yok. Bir sonraki iki alt bölümde bunu somutlaştıracağız.
Bu yüzden neredeyse tüm bilgisayarların aritmetik birimleri 2'nin tümleyeni sistemini kullanır.
\end{bilginot}

\subsection{İşaretli Toplama}
\label{sec:isaretli-toplama}

\begin{tanim}{İşaretli Toplama Kuralı (2'nin Tümleyeni):}
İki işaretli sayı, \textbf{işaret bitleri dahil} sıradan ikili toplama ile toplanır. İşaret
basamağından taşan bir uç elde oluşursa, \textbf{atılır}.
\end{tanim}

\begin{ornek}[Dört Durum, 8 Bit Üzerinde]
Aşağıdaki tabloda negatif sayılar (\emph{$-6$} ve \emph{$-13$}) zaten kendi \textbf{2'nin
tümleyeni} gösterimleriyle verilmiştir; yani $11111010$, $-6$'nın bit dizisidir (bir önceki
alt bölümde gördüğümüz yöntemle hesaplanmıştır), rastgele bir dizi değildir. Yeni başlayan
biri için bu netleştirme önemlidir: tabloya bakarken "bu $1$'lerle dolu dizi nereden çıktı"
diye düşünmeyin. Her zaman önce pozitif sayının ikili karşılığını yazıp sonra (gerekiyorsa)
2'nin tümleyenini alarak elde edilirler.
\begin{center}
\begin{tabular}{lll}
\toprule
İşlem & Toplam & Sonuç \\
\midrule
$(+6)+(+13)$: $00000110+00001101$ & $00010011$ & $+19$ \\
$(-6)+(+13)$: $11111010+00001101$ & $1\,00000111$ (uç elde atılır) & $+7$ \\
$(+6)+(-13)$: $00000110+11110011$ & $11111001$ & $-7$ \\
$(-6)+(-13)$: $11111010+11110011$ & $1\,11101101$ (uç elde atılır) & $-19$ \\
\bottomrule
\end{tabular}
\end{center}
Her dört durumda da kural aynı: işaret dahil topla, taşarsa at. Sonucun negatif mi pozitif
mi olduğu, sonucun kendi işaret bitinden ($0$/$1$) doğrudan okunur.
\end{ornek}

\begin{bilginot}[Taşma (Overflow)]
$n$ bitlik iki sayının toplamı, $n$ bite sığmayan (yani $n{+}1$ bit gerektiren) bir sonuç
üretirse \textbf{taşma} oluşur. Kağıt üzerinde sorun değildir (sınırsız genişletebilirsiniz),
ama donanımda bit genişliği sabittir, taşan sonuç kayba uğrar. Örneğin $4$ bitlik
(aralık $-8$ ile $+7$) bir sistemde $(+5)+(+4)=0101+0100=1001$ işlemini yapalım: sonuç
$1001$, işaret biti $1$ olduğu için \emph{negatif} okunur; $2$'nin tümleyeninde
$1001=-7$. Ama $5+4=9$ pozitif olmalıydı! Sonuç $4$ bite sığmadığı ($9>7$, aralığın üst
sınırını aştığı) için taşma oluştu ve anlamsız bir işaret üretti. Taşmanın devrede nasıl
otomatik tespit edildiğini, toplayıcı devresini kurduğumuzda göreceğiz.
\end{bilginot}

\begin{alistirma}[\kolay]
$4$ bit üzerinde $(+3)+(-5)$ işlemini 2'nin tümleyeni kuralıyla yapınız.
\end{alistirma}

\noindent\textbf{Çözüm:}

\noindent\textbf{Adım 1 (her iki sayıyı da $4$ bitte, işaretli biçimde yazın):} $+3=0011$.
$-5$ için önce $+5=0101$ yazılır, sonra $2$'nin tümleyeni alınır: $1$'in tümleyeni
$1010$, buna $+1$ eklenir, $-5=1011$.

\noindent\textbf{Adım 2 (işaret bitleri dahil sıradan ikili toplama yapın):}
\[
0011+1011=1110
\]

\noindent\textbf{Adım 3 (uç elde var mı kontrol edin):} Sonuç ($1110$) $4$ bite sığıyor,
uç elde \textbf{yok}.

\noindent\textbf{Adım 4 (sonucu yorumlayın):} İşaret biti $1$ olduğu için sonuç negatif;
bu durumda bit dizisi doğrudan büyüklüğü göstermez, kendisi bir $2$'nin tümleyeni
gösterimidir. Gerçek büyüklüğü bulmak için $1110$'nin $2$'nin tümleyenini alırız: $1$'in
tümleyeni $0001$, buna $+1$ eklenir, $0010=2$. Dolayısıyla sonuç $-2$'dir.

Doğrulama: $3-5=-2$. \checkmark

\begin{alistirma}[\orta]
$4$ bit üzerinde $(-3)+(+5)$ işlemini yapınız. Bu kez sonuçta gerçek bir \textbf{uç elde}
oluşacak, atılması gerekecek.
\end{alistirma}

\noindent\textbf{Çözüm:} $-3$'ün 2'nin tümleyeni: $+3=0011$, 1'in tümleyeni $1100$,
$+1=1101$. $+5=0101$.
\[
1101+0101=10010
\]
Bu $5$ bit. $4$ bitlik sistemde bir \textbf{uç elde} ($1$) oluştu, atılır. Kalan
$0010=2$. İşaret biti $0$, yani pozitif: sonuç $+2$. Doğrulama: $-3+5=2$. \checkmark

\subsection{İşaretli Çıkarma}
\label{sec:isaretli-cikarma}

\begin{tanim}{İşaretli Çıkarma Kuralı (2'nin Tümleyeni):}
Çıkanın (çıkarılacak sayının) \textbf{2'nin tümleyenini} (işaret biti dahil) alın ve
eksilene \textbf{toplayın}. İşaret basamağından taşan uç elde atılır.
\end{tanim}

Bu kural, çıkarmayı toplamaya indirger: $(\pm A)-(+B)=(\pm A)+(-B)$ ve
$(\pm A)-(-B)=(\pm A)+(+B)$. Bir pozitif sayıyı negatife çevirmek (ya da tam tersi) zaten
2'nin tümleyenini almaktan ibaret olduğu için, çıkarma her zaman "tümleyen al, topla"
işlemine dönüşür.

\begin{ornek}[$(-6)-(-13)$, 8 Bit Üzerinde]
$-6=11111010$, $-13=11110011$. Çıkanın ($-13$) 2'nin tümleyeni alınır (bu, $+13$'ü verir):
$00001101$. Eksilene ($-6$) eklenir:
\[
11111010+00001101 = 1\,00000111
\]
Uç elde atılır: sonuç $00000111=+7$. Doğrulama: $(-6)-(-13)=7$. \checkmark
\end{ornek}

\begin{ornek}[$(+9)-(+13)$, 8 Bit Üzerinde: Uç Eldesiz Durum]
$+9=00001001$, $+13=00001101$. Çıkanın ($+13$) 2'nin tümleyeni alınır (bu, $-13$'ü verir):
$11110011$. Eksilene ($+9$) eklenir:
\[
00001001+11110011 = 11111100
\]
Bu kez $8$ bite sığıyor, uç elde \textbf{yok}; ama bu, bir önceki bölümdeki işaretsiz
tümleyen çıkarmasından farklı olarak burada bir sorun \textbf{değildir}: işaretli çıkarmada
kural her zaman aynıdır (varsa uç eldeyi at), uç eldenin olup olmaması sonucun işaretini
etkilemez, sonuç zaten kendi işaret bitiyle doğru okunur. $11111100$'ün işaret biti $1$,
yani negatif; 2'nin tümleyeni $00000100=4$, sonuç $-4$. Doğrulama: $9-13=-4$. \checkmark
\end{ornek}

\begin{alistirma}[\kolay]
$(+15)-(+6)$ işlemini $8$ bit üzerinde 2'nin tümleyeni kuralıyla yapınız.
\end{alistirma}

\noindent\textbf{Çözüm:} $+15=00001111$. Çıkanın ($+6$) 2'nin tümleyeni: $+6=00000110$,
1'in tümleyeni $11111001$, $+1=11111010$ ($-6$). Eksilene eklenir:
\[
00001111+11111010 = 1\,00001001
\]
Uç elde atılır: sonuç $00001001=+9$. Doğrulama: $15-6=9$. \checkmark

\begin{alistirma}[\orta]
$(+4)-(+9)$ işlemini $8$ bit üzerinde 2'nin tümleyeni kuralıyla yapınız.
\end{alistirma}

\noindent\textbf{Çözüm:} $+4=00000100$. Çıkanın ($+9$) 2'nin tümleyeni: $+9=00001001$,
1'in tümleyeni $11110110$, $+1=11110111$ ($-9$). Eksilene eklenir:
\[
00000100+11110111 = 11111011
\]
$8$ bite sığıyor, uç elde yok --- önceki örnekte gördüğümüz gibi bu bir sorun değil, kural
aynı şekilde işliyor. İşaret biti $1$, yani negatif: $11111011$'in 2'nin tümleyeni
$00000101=5$, sonuç $-5$. Doğrulama: $4-9=-5$. \checkmark

\begin{bilginot}[Büyük Resim: Tek Bir Devre Yeter]
İşaretli sayılar, işaretsiz sayılarla \textbf{birebir aynı toplama donanımını} kullanır.
Çıkarma bile aynı toplayıcı üzerinden, sadece bir tümleyen alma adımıyla yürür. Bölüm~
\ref{sec:tumleyenler}'in başında sorduğumuz soruya ("neden çıkarmayı toplamaya
çeviriyoruz?") artık tam cevabımız var: hem işaretsiz hem işaretli aritmetik, aynı tek
toplayıcı devre üzerinden, iki ayrı devre kurmaya hiç gerek kalmadan yürüyor. Bu, neredeyse
tüm bilgisayarların aritmetik birimlerinde 2'nin tümleyeni sisteminin evrensel olarak
kullanılmasının temel nedenidir.
\end{bilginot}

\section{IEEE 754 ile Kayan Noktalı Sayı Gösterimi}
\label{sec:ieee754}

Bölüm~\ref{sec:kesirli-donusum}'da, bazı kesirli sayıların ikilide hiçbir zaman
sonlanmadığını görmüş, bu durumun kayan noktalı gösterimle olan ilişkisini ise ileriki bir
bölüme bırakmıştık. Şimdi sırası geldi: o zaman bahsettiğimiz konuya, kayan noktalı sayı
gösterimine, bu bölümde değineceğiz.

Kayan noktalı gösterimin temelinde üç kavram vardır: bir \textbf{işaret}, normalize edilmiş
bir \textbf{anlamlı kısım} (mantissa) ve ölçeği belirleyen bir \textbf{üs}. Bu üç kavramdan
en alışık olmadığımız kısım üstür, çünkü IEEE 754 standardında üs, ham (işaretli) hâliyle
değil, \emph{bias} (kaydırma) adı verilen sabit bir değer eklenerek saklanır: gerçek üse bias
eklenip elde edilen yeni değere \emph{biaslı üs} denir, ve bit dizisinde fiilen saklanan
değer budur, gerçek üs değil. Örneğin tek duyarlıklı formatta bias $127$'dir; gerçek üs
$-1$ ise, saklanan biaslı üs $-1+127=126$ olur. Bu tercihin mühendislik gerekçesini
Bölüm~\ref{sec:ieee754-format}'te göreceğiz; şimdilik terimin ne anlama geldiğini bilmemiz
yeterlidir.

\begin{bilginot}[Üç Terimin Kökeni: Mantissa, Bias, Üs]
Bu üç terimin öyküsü, birbirinden bağımsız üç dilden gelip burada bir araya geliyor.

\textbf{Mantissa}, Etrüsk kökenli bir Latince sözcüktür; aslında "önemsiz bir ek, kefeyi
tamamlamak için eklenen ıvır zıvır" anlamına gelir. 17. yüzyılda logaritma cetvelleri
hazırlanırken, bir logaritmanın tam sayı kısmına ("karakteristik", ondalık noktanın
konumunu belirler) karşılık, ondalık kısmına bu ad verilmişti; çünkü tam sayı kısmının
yanında önemsiz bir "ek" gibi görünüyordu. Oysa asıl anlamlı basamaklar tam da o
kısımdaydı: en önemli parça, en önemsiz görünen adı almış oldu. Bu köken uyuşmazlığı
yüzünden IEEE 754 standardının kendisi bile "mantissa" yerine "significand" (anlamlı
kısım) terimini önerir; çünkü mantissa tarihsel olarak bir logaritmanın ondalık kısmı
içindi, kayan nokta gösterimi ise logaritmik değildir.

\textbf{Bias}, Eski Fransızca "biais" (eğik, çapraz) sözcüğünden İngilizceye geçmiştir;
bilye oyunlarında topun içine kasıtlı asimetrik ağırlık konarak düz değil eğri gitmesini
sağlamak anlamında kullanılıyordu. IEEE 754'teki kullanımı bu köke tam uyar: gerçek üs
değeri bir sabitle "eğilir" (kaydırılır), böylece işaretsiz tam sayı gibi
karşılaştırılabilir hâle gelir.

\textbf{Üs} ise üçünün en şaşırtıcısıdır. Sözcük Arapça "ˀuss" kökünden gelir ve
doğrudan "temel, taban, esas" anlamına gelir (Aramice/Süryanice "Awş" ve Akatça "uşşu"
ile aynı kökten; her ikisi de "temel" demektir); Türkçedeki ilk yazılı kaydı 1680
tarihlidir. Askeri "üs" (karargâh, temel nokta) ile matematikteki "üs" tam olarak aynı
kelimedir. İlginç olan şudur: Türkçe matematik terminolojisinde "taban" zaten ayrı bir
terim olarak $2^3$'teki $2$'yi karşılarken, "üs" kelimesinin kendisi etimolojik olarak
neredeyse tam bu anlamı taşır, yani kuvveti değil "temel/taban" anlamını. Ayrıca "üs"
ile "üst" (yukarı, üzerinde) sıkça karıştırılır, ama bu ikisinin kökeni bile ortak
değildir: "üst" muhtemelen Eski Türkçe "üze" (üzerinde) kökünden gelir, "üs" ise
yukarıda anlatıldığı gibi tamamen ayrı, Arapça kökenli bir sözcüktür. Benzerlik yalnızca
yazılıştan kaynaklanır.
\end{bilginot}

\subsection{Neden Sabit Nokta Yetmez?}
\label{sec:neden-kayan-nokta}

Bölüm~\ref{sec:kesirli-donusum}'da öğrendiğimiz kesirli gösterimde, ikili noktanın konumu
\textbf{sabittir}. Örneğin $8$ biti tam sayı kısmına, $8$ biti kesir kısmına ayırdığımızı
düşünelim: bu düzende hem çok büyük bir sayıyı (milyonlarca gibi) hem çok küçük, hassas bir
kesri (milyonda bir gibi) aynı anda temsil edemeyiz. Tam kısma ayrılan $8$ bit en fazla
$255$'e kadar sayı tutabilir, kesir kısmına ayrılan $8$ bit de belirli bir hassasiyetin
ötesine hiç geçemez. Yani sabit noktalı gösterimde \textbf{aralık ile hassasiyet arasında
sabit, değiştirilemez bir ödünleşim} vardır.

Kayan noktalı gösterim bu sorunu, bilimsel gösterimden ödünç alınan bir fikirle çözer: çok
büyük bir sayıyı $1.234\times10^5$ olarak, çok küçük bir sayıyı da $6.02\times10^{-23}$
olarak yazabildiğimiz gibi, sayıyı doğrudan saklamak yerine \textbf{üç parçaya ayırıp}
saklarız: işaret, normalize edilmiş anlamlı kısım ve üs. Aynı sayıda bit ile, üs değeri
gerektiğinde noktayı çok sola kaydırarak çok küçük sayıları, gerektiğinde çok sağa
kaydırarak çok büyük sayıları temsil etmemizi sağlar. Sabit noktalı gösterimden farklı
olarak, noktanın sayı içindeki gerçek konumu artık tek bir yere sabitlenmemiştir: hangi
basamağın önünde durduğu, her sayı için üs değerine göre değişir, yani noktanın konumu
sayıdan sayıya hareket eder. Gösterimin adı da buradan gelir: nokta tek bir yerde sabit
durmaz, "kayar" (yani İngilizcesiyle "float"lar).

\begin{tanim}{Kayan Noktalı Sayının Genel Formu:}
\[
(-1)^S \times 1.M \times 2^{E-\text{bias}}
\]
Burada $S$ işaret biti ($0$ pozitif, $1$ negatif), $M$ mantissa (anlamlı kısmın noktadan
sonraki basamakları), $E$ saklanan biaslı üs, ve \emph{bias} bir önceki paragrafta
tanımladığımız sabit kaydırma değeridir.
\end{tanim}

\subsection{Normalizasyon: Bedava Bir Bit}
\label{sec:normalizasyon}

Normalizasyon, bir sayıyı $1.xxxx\times2^k$ biçimine getirmek için ikili noktayı
kaydırma işlemidir. Kaydırma yönü tek bir kurala bağlıdır: amaç, noktayı her zaman
sayının \textbf{ilk} (en soldaki) $1$'inin hemen sağına getirmektir; nokta hangi yönde
bu konuma daha yakınsa o yöne kaydırılır.
\begin{itemize}[nosep, leftmargin=*]
    \item Sayının tam kısmında birden fazla basamak varsa (yani sayı $2$'ye eşit ya da
    ondan büyükse), ilk $1$ noktanın \textbf{solundadır}. Nokta, bu $1$'in hemen sağına
    gelene kadar \textbf{sola} kaydırılır. Kaydırılan basamak sayısı, üsse
    \textbf{pozitif} olarak yansır (çünkü sayıyı eski hâline getirmek için sonradan
    $2$'nin o kadar kuvvetiyle çarpmamız gerekir).
    \item Sayı $1$'den küçükse (yani $0.xxxx$ biçimindeyse), ilk $1$ noktanın
    \textbf{sağında}, bir veya daha fazla sıfırdan sonra gelir. Nokta, bu $1$'in hemen
    sağına gelene kadar \textbf{sağa} kaydırılır. Kaydırılan basamak sayısı, üsse
    \textbf{negatif} olarak yansır.
\end{itemize}
Her iki durumda da üs, noktanın kaç basamak ve hangi yönde kaydırıldığının bir kaydıdır;
aşağıdaki örnekte bu iki durumu da (biri sola, biri sağa kaydırma gerektiren) ayrı ayrı
göreceğiz.

\begin{bilginot}[Örtük Baştaki $1$]
İkilide (sıfır hariç) her sayı, tek bir şekilde $1.xxxx\times2^k$ biçiminde normalize
edilebilir; çünkü ikili sistemde ilk anlamlı basamak için tek bir seçenek vardır: ondalıkta
ilk basamak $1$'den $9$'a kadar herhangi bir değer olabilirken, ikilide $0$'dan farklı tek
rakam $1$'dir. Bu evrensel gerçek sayesinde, IEEE 754 formatında baştaki bu $1$ hiçbir zaman
saklanmaz. $M$ (mantissa) alanına yalnızca $1$'den sonraki basamaklar yazılır. IEEE 754
standardının kendisi tek duyarlıklı formatta bu $M$ alanına tam olarak $23$ bit ayırır
(Tablo~\ref{tab:ieee754-format}); standardın doğrudan belirttiği sayı budur, $24$ değil.
$24$, bizim buradan türettiğimiz bir sonuçtur: saklanan $23$ bite, saklanmayan ama
normalizasyon sayesinde her zaman var olduğunu bildiğimiz baştaki $1$'i (yani $1$ bit
daha) eklersek, anlamlı kısmın \emph{gerçek} hassasiyeti $23+1=24$ bit olur. Yani
donanımda yalnızca $23$ bit harcayarak $24$ bitlik bir hassasiyete sahip oluruz: baştaki
$1$ "örtük" (implicit) olduğu için bedavadan kazanılmış bir bit budur. (Aynı mantıkla,
çift duyarlıkta $M$ alanı $52$ bit olduğundan gerçek hassasiyet $52+1=53$ bittir.)
\end{bilginot}

\begin{ornek}[İki Normalizasyon Örneği: Büyük ve Küçük Bir Sayı]
\textbf{Büyük bir sayı:} $(1101.011)_2$'yi normalize edelim. Nokta $3$ basamak sola kaydırılır
(ilk $1$'in hemen sağına):
\[
(1101.011)_2 = 1.101011 \times 2^{3}
\]
Kontrol: $2^3=8$ ile çarpmak, ikili noktayı $3$ basamak sağa kaydırmakla aynı şeydir
(tıpkı onlukta bir sayıyı $10^3$ ile çarpmanın, ondalık noktayı $3$ basamak sağa
kaydırmasına benzer şekilde). $1.101011\times2^3$ işleminde nokta $3$ basamak sağa
kayar:
\[
1.101011 \;\xrightarrow{3\text{ basamak sağa}}\; 1101.011
\]
Bu, normalize etmeden önceki orijinal sayıyla birebir eşleşir. Aynı sonucu ondalık
tabanda çarparak da doğrulayabiliriz:
\[
1.101011_2=1+0.5+0.125+0.03125+0.015625=1.671875
\]
\[
1.671875\times8=13.375
\]
\[
1101.011_2=8+4+1+0.25+0.125=13.375
\]
İki taraf da eşit çıktığından doğrulanmış olur. \checkmark

\textbf{Küçük bir sayı:} $(0.000101)_2$'yi normalize edelim. Nokta $4$ basamak sağa
kaydırılır (ilk $1$'in hemen sağına gelene kadar):
\[
(0.000101)_2 = 1.01 \times 2^{-4}
\]
Kontrol: $1.01\times2^{-4}=1.25\times0.0625=0.078125$, ve $(0.000101)_2=2^{-4}+2^{-6}=
0.0625+0.015625=0.078125$, eşleşiyor.

İki örnek birlikte gösteriyor ki: sayı ne kadar büyükse üs o kadar pozitif, sayı ne kadar
küçükse (sıfıra o kadar yakınsa) üs o kadar negatif olur. Kayan nokta, üssü ayarlayarak
her iki uca da uzanabilir.
\end{ornek}

\subsection{IEEE 754 Formatları ve Bias}
\label{sec:ieee754-format}

\begin{bilginot}[Neden İşaretli Değil, Biaslı Üs?]
Üs için 2'nin tümleyeni yerine biaslı gösterim kullanılmasının mühendislik gerekçesi
şudur: bu sayede iki kayan noktalı sayı, \textbf{sıradan işaretsiz tam sayı
karşılaştırıcısıyla} karşılaştırılabilir. Biaslı gösterimde büyüklük sıralaması korunur.
2'nin tümleyeni kullanılsaydı, negatif üsler (en soldaki bit $1$) pozitif üslerden büyükmüş
gibi görünürdü ve karşılaştırma devresi ek mantık gerektirirdi. Bias değeri, $n$ bitlik üs
alanı için $2^{n-1}-1$ formülüyle seçilir; tek duyarlıkta bu $2^7-1=127$ değerini verir.
\end{bilginot}

\begin{ornek}[Biaslı Üsün Karşılaştırmadaki Avantajı: $e=-2$ ile $e=+3$]
İki kayan noktalı sayının gerçek üsleri $e_1=-2$ ve $e_2=+3$ olsun ($e_2>e_1$, yani ikinci
sayının büyüklük mertebesi daha yüksektir). Tek duyarlıkta $8$ bitlik üs alanı ve
$\text{bias}=127$ kullanarak saklanan biaslı üsleri bulalım:
\[
E_1 = e_1+127 = -2+127 = 125 = (01111101)_2
\]
\[
E_2 = e_2+127 = 3+127 = 130 = (10000010)_2
\]
Bu iki bit dizisi \textbf{sıradan işaretsiz tam sayı} olarak karşılaştırılır: $130>125$,
yani $E_2>E_1$; bu da gerçek üsler arasındaki $e_2>e_1$ ilişkisiyle birebir uyuşur.
\checkmark

Şimdi aynı üsleri, bias \textbf{kullanmadan}, doğrudan $8$ bitlik $2$'nin tümleyeniyle
saklasaydık ne olacağını görelim:
\[
e_1=-2 \;\Rightarrow\; (11111110)_2, \qquad e_2=+3 \;\Rightarrow\; (00000011)_2
\]
Bu iki diziyi işaretsiz tam sayı olarak karşılaştırırsak: $(11111110)_2=254$ ve
$(00000011)_2=3$; $254>3$ görünür, yani $e_1>e_2$ sonucuna varılır. Ama gerçekte
$e_2>e_1$'dir ($3>-2$). İşaretsiz karşılaştırıcı, $2$'nin tümleyeninde yanlış sonuç
verir; bilgi notunda anlatılan sorun tam olarak budur. Bias, üsleri işaretsiz tam sayı
uzayına kaydırarak bu sorunu ortadan kaldırır.
\end{ornek}

\begin{table}[H]
\centering
\begin{tabular}{lccc}
\toprule
Format & İşaret & Üs & Mantissa \\
\midrule
Tek duyarlık (32 bit) & 1 bit & 8 bit (bias $=127$) & 23 bit \\
Çift duyarlık (64 bit) & 1 bit & 11 bit (bias $=1023$) & 52 bit \\
\bottomrule
\end{tabular}
\caption{IEEE 754 Tek ve Çift Duyarlıklı Formatlar}
\label{tab:ieee754-format}
\end{table}

\subsection{Ondalıktan IEEE 754'e Dönüşüm}
\label{sec:ondalik-ieee754}

\begin{ornek}[$(0.6875)_{10}$'un Tek Duyarlıklı IEEE 754 Gösterimi]
Bu sayının ikili karşılığını zaten biliyoruz (Bölüm~\ref{sec:kesirli-donusum}):
$(0.6875)_{10}=(0.1011)_2$.

\textbf{Normalize et:} nokta bir basamak sağa kaydırılır, üssü telafi etmek için $2^{-1}$
çarpanı eklenir:
\[
0.1011 = 1.011\times2^{-1}
\]

\textbf{İşaret:} sayı pozitif olduğu için işaret biti $S=0$ olarak belirlenir.

\textbf{Üs:} gerçek üs $-1$'dir; biaslı üs $-1+127=126$ olarak hesaplanır, bu değer $8$
bitte $(01111110)_2$ şeklinde yazılır.

\textbf{Mantissa:} normalize edilmiş formdaki baştaki $1$'den sonraki basamaklar ($011$)
alınır, $23$ bite tamamlanana kadar sıfırla doldurulur: $01100000000000000000000$.

Üç alan yan yana yazıldığında $32$ bitlik sonuç elde edilir:
\[
\underbrace{0}_{S} \;\; \underbrace{01111110}_{E} \;\; \underbrace{01100000000000000000000}_{M}
\]
\end{ornek}

\begin{alistirma}[\orta]
$(2.5)_{10}$'un tek duyarlıklı IEEE 754 gösterimini bulunuz.
\end{alistirma}

\noindent\textbf{Çözüm:}

\noindent\textbf{Adım 1 (ondalığı ikiliye çevirin):} Tam kısım $2=(10)_2$, kesirli kısım
$0.5=(0.1)_2$; birleştirilince $2.5=(10.1)_2$.

\noindent\textbf{Adım 2 (normalize edin):} Nokta $1$ basamak sola kaydırılır:
$10.1=1.01\times2^{1}$.

\noindent\textbf{Adım 3 (işaret bitini belirleyin):} Sayı pozitif olduğundan $S=0$.

\noindent\textbf{Adım 4 (biaslı üsü hesaplayın):} Gerçek üs $1$'dir; biaslı üs
$1+127=128=(10000000)_2$'dir.

\noindent\textbf{Adım 5 (mantissa alanını doldurun):} Normalize edilmiş anlamlı kısmın
noktadan sonraki basamakları ("$01$"), $23$ bite tamamlanana kadar sağdan sıfırla
doldurulur:
\[
M = 01000000000000000000000
\]

\noindent\textbf{Adım 6 (üç alanı birleştirin):} $S=0$, $E=10000000$, $M$ yukarıdaki
gibidir.
\[
0 \;\; 10000000 \;\; 01000000000000000000000
\]

\begin{alistirma}[\zor]
$(-1.5)_{10}$'un tek duyarlıklı IEEE 754 gösterimini bulunuz.
\end{alistirma}

\noindent\textbf{Çözüm:}

\noindent\textbf{Adım 1 (büyüklüğü ikiliye çevirin):} $|-1.5|=1.5=(1.1)_2$.

\noindent\textbf{Adım 2 (normalize edin):} Sayı zaten $1.1\times2^{0}$ biçimindedir; ek
bir kaydırma gerekmez.

\noindent\textbf{Adım 3 (işaret bitini belirleyin):} Sayı negatif olduğundan $S=1$.

\noindent\textbf{Adım 4 (biaslı üsü hesaplayın):} Gerçek üs $0$'dır; biaslı üs
$0+127=127=(01111111)_2$'dir.

\noindent\textbf{Adım 5 (mantissa alanını doldurun):} Normalize edilmiş anlamlı kısmın
noktadan sonraki tek basamağı ("$1$"), $23$ bite tamamlanana kadar sağdan sıfırla
doldurulur:
\[
M = 10000000000000000000000
\]

\noindent\textbf{Adım 6 (üç alanı birleştirin):} $S=1$, $E=01111111$, $M$ yukarıdaki
gibidir.
\[
1 \;\; 01111111 \;\; 10000000000000000000000
\]

\subsection{IEEE 754'ten Ondalığa Dönüşüm}
\label{sec:ieee754-ondalik}

\begin{ornek}[Bit Dizisinin Çözümlenmesi]
$1\ 10000001\ 01000000000000000000000$ dizisini çözelim.

\textbf{İşaret:} $S=1$ olduğundan sayı negatiftir.

\textbf{Üs:} biaslı üs $(10000001)_2=129$'dur; buradan gerçek üs $129-127=2$ olarak bulunur.

\textbf{Mantissa:} $M=01\ldots0$'dır; örtük baştaki $1$ ile birlikte anlamlı kısım $1.01$
olarak okunur.

Bu üç alan birleştirildiğinde değer şu şekilde hesaplanır:
\[
(-1)^1\times1.01_2\times2^2 = -(101)_2 = -5
\]
\end{ornek}

\begin{alistirma}[\orta]
$0\ 10000010\ 11000000000000000000000$ dizisini çözümleyip ondalık değerini bulunuz.
\end{alistirma}

\noindent\textbf{Çözüm:}

\noindent\textbf{İşaret:} $S=0$ olduğundan sayı pozitiftir.

\noindent\textbf{Üs:} Biaslı üs $(10000010)_2=130$'dur; buradan gerçek üs $130-127=3$
olarak bulunur.

\noindent\textbf{Mantissa:} $M=11\ldots0$'dır; örtük baştaki $1$ ile birlikte anlamlı
kısım $1.11$ olarak okunur.

\noindent Bu üç alan birleştirildiğinde değer şu şekilde hesaplanır:
\[
(-1)^0\times1.11_2\times2^3 = (1110)_2 = 14
\]

\subsection{Özel Değerler ve Kapanış}
\label{sec:ieee754-ozel}

\begin{tanim}{Özel Bit Kalıpları:}
Üs alanının tamamen $0$ veya tamamen $1$ olması özel anlamlar taşır; her durum için tek
duyarlıklı ($32$ bit) bir örnek aşağıda verilmiştir.
\begin{itemize}[nosep, leftmargin=*]
    \item \textbf{Üs tümü $0$, mantissa tümü $0$:} $\pm0$ değerini temsil eder, işaret biti
    hangisi olduğunu belirler (ör. $0\ 00000000\ 0\ldots0$ ifadesi $+0$'ı,
    $1\ 00000000\ 0\ldots0$ ifadesi $-0$'ı verir). İlginçtir: 2'nin tümleyeninde tek bir
    sıfır kodu elde etmek için özenle tasarım yapmıştık (Bölüm~\ref{sec:neden-2-tumleyeni}),
    ama IEEE 754'te $+0$ ile $-0$ yine iki farklı bit kalıbı olarak karşımıza çıkar.
    \item \textbf{Üs tümü $1$, mantissa tümü $0$:} $\pm\infty$ değerini temsil eder (ör.
    $0\ 11111111\ 0\ldots0$ ifadesi $+\infty$'u verir). Örneğin çok büyük iki sayının
    çarpımı, temsil edilebilir aralığı aştığında bu değer üretilir.
    \item \textbf{Üs tümü $1$, mantissa sıfır değil:} NaN (Not a Number) değerini temsil
    eder (ör. $0\ 11111111\ 10\ldots0$), $0/0$ ya da $\infty-\infty$ gibi tanımsız
    işlemlerin sonucudur.
    \item \textbf{Üs tümü $0$, mantissa sıfır değil:} \emph{denormalize} (subnormal) sayıları
    temsil eder (ör. $0\ 00000000\ 0\ldots01$). Bu durumda örtük baştaki $1$ varsayımı
    bırakılır ve değer $(-1)^S\times0.M\times2^{1-\text{bias}}$ formülüyle hesaplanır.
    Amaç, sıfıra çok yakın sayıların aniden değil kademeli olarak küçülmesini sağlamaktır.
\end{itemize}
\end{tanim}

NaN kavramı soyut kalmasın; IEEE 754'ün tanımladığı en tuhaf davranışlardan biri, üç yaygın
programlama dilinde de birebir aynı şekilde karşımıza çıkar: \textbf{NaN, kendisine bile eşit
değildir} ($x \neq x$). Bu, bir hata değil, standardın kendisinin bilinçli bir tasarım
kararıdır. Bu yüzden her dil, eşitlik operatörü yerine ayrı bir \texttt{isnan}/\texttt{isNaN}
fonksiyonu sağlar.

\textbf{C:}
\begin{lstlisting}[language=C]
#include <stdio.h>
#include <math.h>

int main(void) {
    double x = 0.0 / 0.0;                 /* NaN uretir */
    printf("x = %f\n", x);                /* "x = nan" */
    printf("x == x: %d\n", x == x);       /* 0 (yanlis!) */
    printf("isnan(x): %d\n", isnan(x));   /* 1 (dogru yontem) */
    return 0;
}
\end{lstlisting}
\noindent\textbf{Çıktı:}
\begin{lstlisting}[style=cikti]
x = nan
x == x: 0
isnan(x): 1
\end{lstlisting}
\captionof{lstlisting}{C'de NaN üretimi ve $x\neq x$}

\textbf{C++:}
\begin{lstlisting}[language=C++]
#include <iostream>
#include <cmath>
#include <limits>

int main() {
    double x = std::numeric_limits<double>::quiet_NaN();
    std::cout << "x = " << x << "\n";
    std::cout << "x == x: " << (x == x) << "\n";        // false
    std::cout << "isnan(x): " << std::isnan(x) << "\n"; // true
    return 0;
}
\end{lstlisting}
\noindent\textbf{Çıktı:}
\begin{lstlisting}[style=cikti]
x = nan
x == x: 0
isnan(x): 1
\end{lstlisting}
\captionof{lstlisting}{C++'ta NaN üretimi ve $x\neq x$}

\textbf{Java:}
\begin{lstlisting}[language=Java]
public class NanOrnegi {
    public static void main(String[] args) {
        double x = 0.0 / 0.0;                                  // NaN uretir
        System.out.println("x = " + x);                        // "x = NaN"
        System.out.println("x == x: " + (x == x));             // false
        System.out.println("Double.isNaN(x): " + Double.isNaN(x)); // true
    }
}
\end{lstlisting}
\noindent\textbf{Çıktı:}
\begin{lstlisting}[style=cikti]
x = NaN
x == x: false
Double.isNaN(x): true
\end{lstlisting}
\captionof{lstlisting}{Java'da NaN üretimi ve $x\neq x$}

Üçünde de aynı sonuç: $x==x$ karşılaştırması \texttt{false}/$0$ döner, oysa her tür sağduyu
bize bir sayının kendisine eşit olması gerektiğini söyler. Bunun nedeni, IEEE 754
standardının NaN karşılaştırmalarını \emph{tanımsız} kabul etmesi ve tüm karşılaştırma
operatörlerinin ($==$, $<$, $>$ vb.) NaN içeren her işlemde \texttt{false} döndürmesini
zorunlu kılmasıdır. $\neq$ operatörü bu kuralın tek istisnasıdır ve NaN için her zaman
\texttt{true} döner. Bu yüzden bir değişkenin NaN olup olmadığını kontrol etmek için asla
\texttt{x == x} gibi bir eşitlik testi değil, dilin sağladığı özel fonksiyon
(\texttt{isnan}/\texttt{std::isnan}/\texttt{Double.isNaN}) kullanılmalıdır.

\begin{bilginot}[Artık Biliyoruz: $0.1+0.2\neq0.3$]
Daha önce (Bölüm~\ref{sec:kesirli-donusum}) $(0.1)_{10}$'un ikilide hiç sonlanmadığını
görmüştük. IEEE 754'te mantissa \textbf{sonlu} ($23$ ya da $52$ bit) olduğundan, sonlanmayan
bu sayı en yakın temsil edilebilir değere \textbf{yuvarlanır}; gerçek $0.1$'den son derece
küçük ama sıfır olmayan bir farkla. $0.1$ ve $0.2$ ayrı ayrı yuvarlandıktan sonra toplanınca,
bu küçük yuvarlama hataları birikir ve sonuç tam olarak $0.3$'ün temsiline eşit çıkmaz. Bu,
bir yazılım hatası değil, sonlu bit genişliğinin kaçınılmaz bir sonucudur.
\end{bilginot}

\section{İkili Kodlama Sistemleri}
\label{sec:ikili-kodlar}

Şimdiye kadar Bölüm~\ref{chap:sayi-sistemleri} boyunca hep \emph{sayıları} ikiliye çevirmekle
uğraştık. Ama bir bilgisayarın işlediği bilginin çoğu aslında sayısal bir büyüklük değildir:
harfler, semboller, ya da "hangi çalışan", "hangi renk" gibi kategorik kimlikler de ikili
düzende temsil edilmelidir. Bu bölümde, sayısal değeri değil bir \emph{kimliği} temsil eden
bit dizilerini, yani \textbf{kodları}, inceleyeceğiz.

\begin{tanim}{İkili Kod:}
Bir kümede $N$ farklı eleman varsa, bu elemanları ayırt eden bir $n$-bit \emph{ikili kod},
her elemana $2^n$ olası bit kombinasyonundan benzersiz birini atar. Ayırt etmek için gereken
minimum bit sayısı $n=\lceil\log_2 N\rceil$'dir, çünkü $n$ bit tam olarak $2^n$ farklı
kombinasyon üretir. Ancak \textbf{maksimum} bit sayısına dair bir sınır yoktur: $N$ elemanı
kodlamak için gerekenden fazla bit de kullanılabilir, fazladan kombinasyonlar basitçe
kullanılmadan bırakılır.
\end{tanim}

Örneğin, 10 ondalık rakamı (0-9) ayırt etmek için teorik minimum $n=4$ bittir, çünkü
$2^3=8<10\le2^4=16$. Bu durumda 16 olası kombinasyondan yalnızca 10 tanesi kullanılır;
geri kalan 6 kombinasyon ($1010$'dan $1111$'e kadar) tanımsız/kullanılmaz kalır. Birazdan
göreceğimiz BCD kodu tam olarak bu stratejiyi izler.

\subsection{BCD (Binary-Coded Decimal) Kodu}
\label{sec:bcd}

En yaygın kullanılan ondalık kod, her ondalık rakamı \textbf{ayrı ayrı} 4 bitlik ikili
karşılığına çevirip bu grupları yan yana dizen \textbf{BCD} (Binary-Coded Decimal,
ikili-kodlanmış ondalık) kodudur. BCD'ye bazen \textbf{8421 kodu} da denir, çünkü her bitin
ağırlığı sırasıyla $8,4,2,1$'dir; yani her 4-bitlik grup, doğrudan o rakamın kendi ikili
karşılığıdır.

\begin{tanim}{BCD Kodu:}
$k$ ondalık haneli bir sayının BCD karşılığı $4k$ bit gerektirir; her hane, $0000$'dan
$1001$'e kadarki (yani 0'dan 9'a kadarki) 4-bit ikili karşılığıyla ayrı ayrı kodlanır.
$1010$'dan $1111$'e kadarki 6 kombinasyon BCD'de \textbf{geçersizdir}, hiçbir anlamı yoktur.
\end{tanim}

\textbf{Önemli:} BCD, bir sayının ikili karşılığıyla \textbf{aynı şey değildir}. Bunu bir
örnekle netleştirelim.

\begin{ornek}[$(185)_{10}$'in BCD ve İkili Karşılıklarının Karşılaştırılması]
Gerçek ikili karşılık, sayının bütününü tek seferde ikiliye çevirerek bulunur:
\[
(185)_{10} = (10111001)_2 \quad \text{(8 bit yeterli)}
\]
BCD karşılık ise her haneyi \textbf{ayrı ayrı} kodlayıp yan yana dizerek bulunur:
\[
1\to0001,\quad 8\to1000,\quad 5\to0101 \quad\Rightarrow\quad (185)_{10}=(0001\ 1000\ 0101)_{BCD}
\]
BCD karşılık 12 bit gerektirirken gerçek ikili karşılık yalnızca 8 bit yeterlidir; aynı
sayı için iki tamamen farklı bit deseni.
\end{ornek}

BCD, gerçek ikiliden daha fazla bit harcadığı halde neden kullanılır? Çünkü BCD bir
\textbf{hız ya da verimlilik} kodu değil, bir \textbf{insan-makine arayüz kolaylığı} kodudur.
Bilgisayarın girdi/çıktısını üreten insan ondalık sistemle düşünür; bir ekrana "185"
yazdırmak ya da klavyeden "1", "8", "5" tuşlarını okumak, sayının tamamını ikiliye çevirip
geri çevirmektense, her hanenin kendi 4-bit kutusunda ayrı ayrı durduğu bir temsille çok
daha doğrudandır.

\begin{alistirma}[\kolay]
$(84)_{10}$ sayısının BCD karşılığını bulunuz.
\end{alistirma}

\noindent\textbf{Çözüm:} $8\to1000$, $4\to0100$, yan yana dizilince $(84)_{10}=(1000\
0100)_{BCD}$.

\subsection{BCD Toplama}
\label{sec:bcd-toplama}

BCD'nin en can sıkıcı yönü burada ortaya çıkar: düz ikili toplama kuralları BCD'de doğrudan
işlemez. İki BCD hanesi (ve olası bir önceki elde) toplandığında, sonuç en fazla $9+9+1=19$
olabilir; bu iki 4-bit grubu sıradan ikili sayılarmış gibi toplarsak, donanım bize doğru
\emph{ikili} toplamı verir ($0$'dan $19$'a kadar). Ama BCD'de yalnızca $0000$-$1001$ (0-9)
geçerlidir. Toplam $1010$ (10) veya üzerindeyse, ya da 4. bitten bir elde taştıysa, elde
edilen sonuç BCD açısından geçersizdir.

\begin{tanim}{BCD Toplama Kuralı (+6 Düzeltmesi):}
İki BCD hanesi (ve varsa önceki elde) sıradan ikili sayılarmış gibi toplanır. Elde edilen
4-bitlik ikili toplam:
\begin{itemize}[nosep, leftmargin=*]
    \item $\le1001$ (9) ise ve 4. bitten bir elde taşmadıysa, sonuç zaten geçerli bir BCD
    hanesidir, düzeltme gerekmez.
    \item $\ge1010$ (10) ise \textbf{veya} 4. bitten bir elde taştıysa, kalan 4 bite $0110$
    (6) eklenir; bu hem doğru BCD hanesini üretir hem de (gerekmiyorsa) bir sonraki haneye
    giden elde bitini doğurur.
\end{itemize}
Düzeltmenin neden tam $6$ olduğu, 4 bitin sağladığı $16$ kombinasyondan yalnızca $10$'unun
BCD'de geçerli olmasından kaynaklanır: $16-10=6$. Toplam geçersiz bölgeye ($10$-$15$)
girdiğinde, $6$ eklemek onu bu altı kullanılmayan kodun üzerinden atlatarak doğru hane ve
elde bitini üretir.
\end{tanim}

\begin{ornek}[Düzeltme Gerekmeyen Toplam: $4+5$]
\[
\begin{array}{r}
0100 \\
+\;0101 \\
\hline
1001
\end{array}
\]
İkili toplam $1001$ (9) $\le1001$ olduğundan zaten geçerli bir BCD hanesidir, düzeltme
gerekmez.
\end{ornek}

\begin{ornek}[Düzeltme Gereken Toplam: $4+8$]
\[
\begin{array}{r}
0100 \\
+\;1000 \\
\hline
1100
\end{array}
\]
İkili toplam $1100$ (12) $\ge1010$ olduğundan geçersizdir; $0110$ (6) eklenir:
\[
\begin{array}{r}
1100 \\
+\;0110 \\
\hline
1\,0010
\end{array}
\]
Sonuç: hane $0010$ (2), elde $1$ --- yani "12" (elde$=1$, hane$=2$). Doğru!
\end{ornek}

\begin{ornek}[Taşan Elde ile Düzeltme: $8+9$]
\[
\begin{array}{r}
1000 \\
+\;1001 \\
\hline
1\,0001
\end{array}
\]
Burada zaten 4. bitten bir elde taşmış durumdadır. Kalan 4 bite yine $0110$ eklenir:
\[
\begin{array}{r}
0001 \\
+\;0110 \\
\hline
0111
\end{array}
\]
Elde zaten mevcuttu, kalan hane $0111$ (7) --- sonuç "17" (elde$=1$, hane$=7$). Doğru!
\end{ornek}

Bu tek haneli kural, çok haneli sayılara şöyle genişler: hanelere sağdan sola, birer birer
uygulanır ve her hanenin ürettiği elde bir sonraki (daha soldaki) hanenin toplamına dahil
edilir, tıpkı okulda öğrendiğimiz elle toplamada olduğu gibi.

\begin{ornek}[Çok Haneli BCD Toplama: $(184)_{BCD}+(576)_{BCD}$]
Sağdan sola, hane hane:
\begin{center}
\small
\begin{tabular}{@{}lp{4.2cm}p{2.6cm}p{2.4cm}@{}}
\toprule
Hane & İkili Toplam (elde dahil) & Düzeltme & Sonuç / Elde \\
\midrule
Birler & $0100+0110=1010$ & $\ge1010$, $+0110$ & $0000$ / $1$ \\
Onlar & $1000+0111+0001=1\,0000$ & elde taştı, $+0110$ & $0110$ / $1$ \\
Yüzler & $0001+0101+0001=0111$ & $\le1001$, gerekmez & $0111$ / $0$ \\
\bottomrule
\end{tabular}
\end{center}
Elde zincirinin izlenmesi önemlidir: birler hanesinin ürettiği elde ($1$), onlar hanesinin
toplamına üçüncü terim olarak girmiştir ($1000+0111+0001$); onlar hanesinin ürettiği elde
($1$) de aynı şekilde yüzler hanesinin toplamına girmiştir ($0001+0101+0001$). Yüzler
hanesinde üretilen elde $0$ olduğundan zincir burada sona erer. Eğer $1$ olsaydı, dördüncü
bir (binler) haneye taşınacaktı. Sağdan sola haneler birleştirilince: $(0111\ 0110\
0000)_{BCD}=(760)_{10}$. Doğrulama: $184+576=760$. \checkmark
\end{ornek}

\begin{alistirma}[\kolay]
$4+6$ toplamının BCD karşılığını (gerekiyorsa düzeltme dahil) bulunuz.
\end{alistirma}

\noindent\textbf{Çözüm:} $0100+0110=1010$ (10), $\ge1010$ olduğundan $0110$ eklenir:
$1010+0110=1\,0000$. Sonuç: hane $0000$ (0), elde $1$; yani BCD toplamı
$(1\,0000)_{BCD}$, ondalıkta $10$.

\subsection{Diğer Ondalık Kodlar}
\label{sec:diger-ondalik-kodlar}

BCD, dört bitle on rakamı kodlamanın \textbf{tek} yolu değildir; sadece en yaygın
olanıdır. Aynı 4 bitten 10 farklı kombinasyon seçmenin başka yolları da vardır, ve her seçim
farklı bir mühendislik avantajı sağlar.

\begin{tanim}{Ağırlıklı Kod:}
BCD'de her bitin sabit bir \emph{ağırlığı} vardır ($8,4,2,1$) ve bir hanenin değeri, $1$
olan bitlerin ağırlıklarını toplayarak bulunur. Aynı fikir farklı ağırlık kümeleriyle de
uygulanabilir; buna \textbf{ağırlıklı kod} denir. Ağırlıkların hepsinin pozitif olması da
gerekmez.
\end{tanim}

\begin{ornek}[2421 Kodu: Aynı Rakam İçin İki Geçerli Kodlama]
2421 kodunda ağırlıklar $2,4,2,1$'dir; bir bit kalıbının değeri, $1$ olan bitlerin
ağırlıklarını toplayarak bulunur.

\noindent\textbf{Birinci kalıp:} $0100$ bit kalıbını değerlendirelim:
\[
2\cdot0+4\cdot1+2\cdot0+1\cdot0=4
\]

\noindent\textbf{İkinci kalıp:} Şimdi $1010$'u değerlendirelim:
\[
2\cdot1+4\cdot0+2\cdot1+1\cdot0=4
\]

\textbf{Aynı değer, iki farklı bit kalıbı!} Rakam $4$, hem $0100$ hem de $1010$ ile
kodlanabilir. 2421 kodunda bazı rakamlar için birden fazla geçerli kodlama vardır; bu seçim
özgürlüğü rastgele değildir. Birazdan göreceğimiz \emph{öz-tümleyen} özelliğini sağlamak için
kasıtlı olarak kullanılır.
\end{ornek}

Ağırlıklı olmayan bir kod da mümkündür: \textbf{Excess-3} (fazla-3) kodunda hiçbir ağırlık
yorumu yoktur, kural yalnızca şudur: \emph{her rakamın düz ikili karşılığına $3$ ekleyin.}
Örneğin rakam $5$: $0101+0011=1000$.

Aşağıdaki tablo, dört farklı ondalık kodu karşılaştırmalı olarak gösterir.

\begin{table}[H]
\centering
\renewcommand{\arraystretch}{1.3}
\begin{tabular}{c@{\hspace{1.4em}}c@{\hspace{1.4em}}c@{\hspace{1.4em}}c@{\hspace{1.4em}}c}
\toprule
Rakam & BCD ($8421$) & $2421$ & Excess-3 & $8,4,-2,-1$ \\
\midrule
0 & 0000 & 0000 & 0011 & 0000 \\
1 & 0001 & 0001 & 0100 & 0111 \\
2 & 0010 & 0010 & 0101 & 0110 \\
3 & 0011 & 0011 & 0110 & 0101 \\
4 & 0100 & 0100 & 0111 & 0100 \\
5 & 0101 & 1011 & 1000 & 1011 \\
6 & 0110 & 1100 & 1001 & 1010 \\
7 & 0111 & 1101 & 1010 & 1001 \\
8 & 1000 & 1110 & 1011 & 1000 \\
9 & 1001 & 1111 & 1100 & 1111 \\
\bottomrule
\end{tabular}
\caption{Dört Farklı Ondalık Kodun Karşılaştırılması}
\label{tab:ondalik-kodlar}
\end{table}

Bu özelliğin neye yaradığını anlamak için önce nereden geldiğini hatırlayalım:
Bölüm~\ref{sec:tumleyenler}'de çıkarmayı toplamaya çevirmek için tümleyen almıştık. Aynı
fikri şimdi \emph{ondalık} sayılar için uygulayacağız; bu da her seferinde bir rakamın
$9$'un tümleyenini ($9-d$) hesaplamayı gerektirir. Bunu donanımda hesaplamanın iki yolu var:
\begin{itemize}[nosep, leftmargin=*]
    \item \textbf{Aritmetik yoldan:} Girdisi $d$, çıktısı $9-d$ olan gerçek bir çıkarma
    devresi (subtractor) kurulur. Bu yol, kapı sayısı bakımından pahalı ve gecikme
    bakımından yavaştır.
    \item \textbf{Kod seçimiyle:} Rakamlar öyle bir kodla saklanırsa ki, kodun bitlerini
    teker teker \textbf{tersine çevirmek} ($1\to0$, $0\to1$) otomatik olarak $9-d$'nin
    kodunu verir; bu durumda çıkarma devresine hiç gerek kalmaz, yalnızca $4$ paralel NOT
    kapısı yeterli olur.
\end{itemize}

Burada terminolojiyi netleştirmek gerekir: bir \textbf{kodlama şeması} (Excess-3, BCD gibi),
her rakama bir \textbf{bit kalıbı} atayan kuraldır. Aşağıdaki tanım bu iki kavramı ayrı
tutarak ifade edilmiştir.

\begin{tanim}{Öz-Tümleyen (Self-Complementing) Kod:}
Bir kodlama şemasında her rakama bir bit kalıbı atanır. Bu şema, herhangi bir rakamın bit
kalıbındaki tüm bitler tersine çevrildiğinde ($1\to0$, $0\to1$) ortaya çıkan yeni bit
kalıbı, tam olarak $9-d$ rakamına atanmış bit kalıbıyla aynı oluyorsa \emph{öz-tümleyen}
olarak adlandırılır.
\end{tanim}

\begin{ornek}[Tanımın Somut Bir Kontrolü: 2421 Kodunda Rakam $4$]
Tanımı, elimizde zaten olan Tablo~\ref{tab:ondalik-kodlar}'daki 2421 kodu üzerinden tek bir
rakamla sınayalım. Rakam $4$'ün 2421 kodundaki bit kalıbı $0100$'dür. Bitlerini tek tek
tersine çevirelim:
\[
0100 \quad\Rightarrow\quad 1011
\]
Şimdi $9-4=5$'in 2421 kodundaki bit kalıbına bakalım: yine Tablo~\ref{tab:ondalik-kodlar}'a
göre $1011$'dir. İki sonuç birebir aynı ($1011=1011$); yani rakam $4$'ün bit kalıbını
çevirmek, bize doğrudan rakam $5$'in bit kalıbını verdi. 2421 kodu, en azından bu rakam
için, tanımdaki öz-tümleyen koşulunu sağlıyor.
\end{ornek}

Tek bir rakam için doğrulamak fikri somutlaştırır, ama genel olarak geçerli olduğunu
kanıtlamaz. Bunu Excess-3 üzerinden, adım adım görelim. Ama önce, bu kodun
tanımında ilk bakışta kafa karıştırabilecek bir noktayı vurgulayalım: Excess-3 kuralı,
\emph{"rakam $d$'nin bit kalıbı, $d+3$ sayısının düz ikili karşılığıdır"} der. Yani bir bit
kalıbının \textbf{sayısal değeri} ile o kalıbın \textbf{temsil ettiği rakam}, aralarında
sabit bir $+3$ fark olacak şekilde, birbirinden farklıdır. Kodun adındaki "excess"
("fazlalık") kelimesi tam olarak bu farkı anlatır.

\begin{ornek}[Excess-3'ün Öz-Tümleyen Özelliği]
Amacımız şunu göstermek: rakam $3$'ün Excess-3 bit kalıbının bitlerini tersine çevirdiğimizde
elde ettiğimiz yeni kalıbın, rakam $9-3=6$'nın Excess-3 bit kalıbıyla birebir aynı çıktığını
görmek. Bunu dört adımda yapalım.

\noindent\textbf{Adım 1 (rakam $3$'ün bit kalıbını bulalım):} Rakam $3$'ün düz ikili karşılığı
$0011$'dir. Excess-3 kuralı gereği buna $3$ eklenir (yani $0011+0011$):
\[
0011+0011=0110
\]
Burada dikkat edilmesi gereken can alıcı nokta şudur: $0110$'nın \textbf{sayısal değeri}
$6$'dır ($4+2=6$). Ama bu bit kalıbı, rakam $6$'yı değil, \textbf{rakam $3$'ü} temsil etmek
üzere atanmıştır, çünkü kuralımız "$d+3$" idi ve $3+3=6$. Yani rakam $3$'ün bit kalıbı,
sayısal değeri $6$ olan $0110$'dır; bu bir çelişki değil, kodun tanımının doğal bir
sonucudur.

\noindent\textbf{Adım 2 (aynı kuralı rakam $6$'ya uygulayalım):} Rakam $6$'nın düz ikili
karşılığı $0110$'dır. Buna $3$ eklenir: $0110+0011=1001$. $1001$'in sayısal değeri $9$'dur,
ama bu bit kalıbı \textbf{rakam $6$'yı} temsil eder (yine $+3$ farkla: $6+3=9$).

\noindent\textbf{Adım 3 (Adım 1'in kalıbının bitlerini çevirelim, hiç toplama yapmadan):}
Rakam $3$'ün bit kalıbı olan $0110$'ı ele alalım ve her bitini tek tek tersine çevirelim:
\[
\underbrace{0}_{\to1}\ \underbrace{1}_{\to0}\ \underbrace{1}_{\to0}\ \underbrace{0}_{\to1}
\quad\Rightarrow\quad 1001
\]

\noindent\textbf{Adım 4 (sonuçları karşılaştıralım):} Adım 2'de toplama yaparak bulduğumuz
rakam $6$'nın bit kalıbı ($1001$) ile Adım 3'te hiç toplama yapmadan, sadece rakam $3$'ün
bit kalıbının bitlerini çevirerek bulduğumuz sonuç ($1001$) birebir aynıdır. Yani rakam
$3$'ün bit kalıbının bitlerini çevirmek, bize doğrudan rakam $9-3=6$'nın bit kalıbını verdi,
hiçbir toplama veya çıkarma yapmadan.
\end{ornek}

Bu bir rastlantı değildir; "$3$" sayısının seçimi de rastgele değildir. Kanıtlayalım.
Bir $4$-bitlik sayının bitlerini tek tek çevirmek, matematiksel olarak $15-x$ işlemine
eşittir (çünkü $4$ bit $0$'dan $15$'e kadar sayar, ve her biti çevirmek sayıyı $15$'e göre
"aynalar": $0000\to1111$ yani $0\to15$, $0001\to1110$ yani $1\to14$, ve genel olarak
$x\to15-x$). Rakam $d$'nin bit kalıbının sayısal değeri $d+3$ olduğundan, bu bit kalıbının
bitlerini çevirmek bize $15-(d+3)=12-d$ değerinde bir bit kalıbı verir. Öte yandan, bizim
aslında elde etmek istediğimiz şey, rakam $9-d$'ye atanmış bit kalıbıdır; bu kalıbın
sayısal değeri de $(9-d)+3=12-d$'dir. \textbf{İkisi de $12-d$; birebir eşit!} Bu sabit
ekleme miktarına (Excess-3'te $3$, BCD'de $0$) kodun \textbf{ofseti} denir. Ofset, bir
bit kalıbının sayısal değeri ile o kalıbın temsil ettiği rakam arasındaki sabit farktır.
İşte Excess-3'ün ofseti olan "$3$", tam olarak yukarıdaki eşitliği ($15-3=9+3$) sağlayacak
şekilde seçilmiştir; rastgele bir tercih değildir.

Bu özellik neden önemlidir? Çünkü ondalık çıkarmada (bir sonraki alt bölümde göreceğimiz
gibi) bir tümleyen alma işlemi normalde ayrı bir çıkarma devresi gerektirir; ama
öz-tümleyen bir kodda bu işlem tek bir NOT kapısı katmanına indirgenir, ekstra donanıma hiç
gerek kalmaz.

\begin{alistirma}[\kolay]
BCD (\emph{8421}) kodunun öz-tümleyen olmadığını, yukarıdaki kanıtla aynı yöntemi
kullanarak gösteriniz.
\end{alistirma}

\noindent\textbf{Çözüm:}

\noindent\textbf{Adım 1 (BCD'nin ofsetini belirleyin):} BCD'nin ofseti $0$'dır, yani bit
kalıbının sayısal değeri rakamın kendisine eşittir ($d+0=d$). Excess-3'teki "$+3$" yerine
burada "$+0$" vardır.

\noindent\textbf{Adım 2 (bit çevirmenin verdiği değeri hesaplayın):} Bit çevirme yine
$15-x$ işlemine eşittir; BCD kodunu çevirirsek $15-d$ elde ederiz.

\noindent\textbf{Adım 3 (elde etmek istediğimiz değerle karşılaştırın):} Bizim istediğimiz,
$9-d$'nin BCD kodu, yani sadece $9-d$'nin kendisidir (ofset $0$ olduğu için). Öz-tümleyen
olması için $15-d=9-d$ eşitliğinin sağlanması gerekirdi; bu ise $d$'den bağımsız olarak
asla doğru olamaz ($15\neq9$).

\noindent\textbf{Adım 4 (somut bir örnekle doğrulayın):} $0$'ın BCD kodu $0000$'dır;
bitleri çevirirsek $1111$ (15) elde ederiz. Ama $9-0=9$'un BCD kodu $1001$'dir;
$1111\neq1001$.

Demek ki BCD'nin öz-tümleyen olmaması tek bir rakama özgü bir tesadüf değil, BCD'nin
ofsetinin ($0$) Excess-3'ün gerektirdiği değere ($3$) eşit olmamasından kaynaklanan
\textbf{yapısal bir imkansızlıktır}.

Son olarak, yukarıdaki tablonun son sütunu bir \textbf{negatif ağırlıklı kod}
örneğidir: $8,4,-2,-1$ kodunda ağırlıkların hepsi pozitif olmak zorunda değildir. Örneğin
$0110$ bu kodda $8\cdot0+4\cdot1+(-2)\cdot1+(-1)\cdot0=2$ olarak yorumlanır --- ve bu kod da
(2421 gibi) öz-tümleyendir.

\subsection{BCD'de İşaretli Ondalık Aritmetik (10'un Tümleyeni)}
\label{sec:bcd-isaretli-aritmetik}

Bölüm~\ref{sec:isaretli-sayilar}'de ikili sayılarda negatifleri saklamak için $2$'nin
tümleyenini kullanmış, çıkarmayı toplamaya çevirmiştik. Aynı fikri şimdi BCD ile saklanan
ondalık sayılara uygulayacağız; bu sefer kullanılan araç \textbf{$10$'un tümleyeni}dir.

\begin{tanim}{İşaret Hanesi:}
Bir işaretli BCD sayısında işaret, ayrı bir bit değil, ayrı bir \textbf{hane} olarak
saklanır: bu hane pozitif sayılar için $0000$, negatif sayılar için $1001$ (rakam $9$'un
BCD kodu) değerini alır.
\end{tanim}

Neden özellikle $9$ seçilmiştir? Çünkü $10$'un tümleyenini alma işlemi, negatif bir sayının
en soluna hep $9$ hanesi(leri) dizer. Bu, $2$'nin tümleyeninde negatif sayıların en soluna
$1$ bitleri dizilmesinin (işaret genişletme) ondalıktaki karşılığıdır.

\begin{tanim}{Bir BCD Sayısının $10$'un Tümleyenini Bulma:}
İki adımlı bir işlemdir:
\begin{enumerate}[nosep, leftmargin=*]
    \item Her hanenin $9$'un tümleyeni alınır (her rakam $d$ için $9-d$ hesaplanır); buna
    \emph{$9$'un tümleyeni} denir.
    \item Sonucun en sağdaki hanesine $1$ eklenir.
\end{enumerate}
\end{tanim}

\begin{bilginot}[Öz-Tümleyen Kodlarla Bağlantı]
Bir önceki alt bölümde BCD'nin öz-tümleyen olmadığını kanıtlamıştık. Bunun donanımsal
bedeli burada ortaya çıkar: yukarıdaki $1$. adımda her hanenin $9$'un tümleyenini bulmak
gerçek bir çıkarma gerektirir. Rakamlar öz-tümleyen bir kodla (Excess-3 gibi) saklansaydı,
bu adım yalnızca bit çevirmeyle (NOT kapılarıyla) yapılabilirdi.
\end{bilginot}

\begin{tanim}{İşaretli Ondalık Toplama Kuralı:}
İşaret hanesi dahil tüm haneler, BCD toplama kuralıyla (gerektiğinde $+6$ düzeltmesiyle)
toplanır. En soldaki haneden (işaret hanesinden) taşan son elde atılır, tıpkı $2$'nin
tümleyeninde işaret bitinden taşan eldenin atılması gibi.
\end{tanim}

\begin{ornek}[$(+375)+(-240)$]
Önce $-240$'ı $10$'un tümleyeniyle temsil edelim. $240$'ın $9$'un tümleyeni: $9-2=7$,
$9-4=5$, $9-0=9$, yani $759$. En sağdaki haneye $1$ eklenir: $760$. İşaret hanesi $9$
olduğundan, $-240\to9760$.

Şimdi $0375+9760$ toplamını hane hane (sağdan sola) yapalım. Hatırlatma: bir hanenin
toplamı $9$'u aşarsa (ya da elde üretirse), Bölüm~\ref{sec:bcd-toplama}'da gördüğümüz
$+6$ BCD düzeltmesi uygulanır:

\begin{center}
\small
\begin{tabular}{@{}lp{3.6cm}p{2.6cm}p{2.8cm}@{}}
\toprule
Hane & Toplam & Düzeltme ($+6$ BCD kuralı) & Sonuç / Elde \\
\midrule
Birler & $5+0=5$ & gerekmez & $5$ / $0$ \\
Onlar & $7+6=13$ & $\ge10$: $+6$ eklenir & $3$ / $1$ \\
Yüzler & $3+7+1=11$ & $\ge10$: $+6$ eklenir & $1$ / $1$ \\
İşaret & $0+9+1=10$ & $\ge10$: $+6$ eklenir & $0$ / $1$ (atılır) \\
\bottomrule
\end{tabular}
\end{center}

Son elde (işaret hanesinden taşan) atılır. Sonuç haneleri (işaret, yüzler, onlar, birler)
$0,1,3,5$'tir, yani $+135$. Doğrulama: $375-240=135$.
\end{ornek}

\begin{alistirma}[\orta]
İki parçada çözünüz: (a) $370+(-250)$, (b) $250+(-370)$. İkinci kısımda sonucun işaret
hanesi negatif çıkacaktır; gerçek büyüklüğü bulmak için sonucun büyüklük hanelerinin
$10$'un tümleyenini alınız.
\end{alistirma}

\noindent\textbf{Çözüm (a):} $-250$'nin $10$'un tümleyeni: $9$'un tümleyeni $9-2=7$,
$9-5=4$, $9-0=9$, yani $749$; en sağa $1$ eklenir: $750$; işaret hanesi $9$: $-250\to9750$.

\begin{center}
\small
\begin{tabular}{@{}lp{3.6cm}p{2.6cm}p{2.8cm}@{}}
\toprule
Hane & Toplam & Düzeltme ($+6$ BCD kuralı) & Sonuç / Elde \\
\midrule
Birler & $0+0=0$ & gerekmez & $0$ / $0$ \\
Onlar & $7+5=12$ & $\ge10$: $+6$ eklenir & $2$ / $1$ \\
Yüzler & $3+7+1=11$ & $\ge10$: $+6$ eklenir & $1$ / $1$ \\
İşaret & $0+9+1=10$ & $\ge10$: $+6$ eklenir & $0$ / $1$ (atılır) \\
\bottomrule
\end{tabular}
\end{center}

Sonuç haneleri $0,1,2,0$'dır, yani $+120$. Doğrulama: $370-250=120$.

\noindent\textbf{Çözüm (b):} $-370$'in $10$'un tümleyeni: $9$'un tümleyeni $9-3=6$,
$9-7=2$, $9-0=9$, yani $629$; en sağa $1$ eklenir: $630$; işaret hanesi $9$: $-370\to9630$.

\begin{center}
\small
\begin{tabular}{@{}lp{3.6cm}p{2.6cm}p{2.8cm}@{}}
\toprule
Hane & Toplam & Düzeltme ($+6$ BCD kuralı) & Sonuç / Elde \\
\midrule
Birler & $0+0=0$ & gerekmez & $0$ / $0$ \\
Onlar & $5+3=8$ & gerekmez & $8$ / $0$ \\
Yüzler & $2+6=8$ & gerekmez & $8$ / $0$ \\
İşaret & $0+9=9$ & gerekmez & $9$ / $0$ \\
\bottomrule
\end{tabular}
\end{center}

Sonuç haneleri $9,8,8,0$'dır. İşaret hanesi $9$ olduğundan sonuç negatiftir; gerçek
büyüklüğü bulmak için büyüklük hanelerinin ($880$) $10$'un tümleyenini alalım: $9$'un
tümleyeni $9-8=1$, $9-8=1$, $9-0=9$, yani $119$; en sağa $1$ eklenir: $120$. Sonuç
$-120$'dir. Doğrulama: $250-370=-120$.

Ondalık çıkarma ($M-N$) da ikilideki gibi yapılır: $N$'nin (çıkanın) $10$'un tümleyeni
alınıp $M$'ye eklenir.

\subsection{Gray Kodu}
\label{sec:gray-kodu}

Düz ikili sayma sırasında, ardışık iki sayı arasında \textbf{birden fazla bit aynı anda}
değişebilir. Örneğin $7\to8$ geçişi $0111\to1000$'dir: dört bitin dördü de aynı anda
değişir.

Bu, kağıt üzerinde bir sorun değildir. Gerçek donanımda ise ciddi bir soruna yol açabilir:
farklı bitlerin mantık kapılarından geçip yeni değerlerine ulaşması tam olarak aynı anda
gerçekleşmez. Üretim toleransları, kapı gecikmeleri ve tel uzunlukları gibi nedenlerle,
bitler arasında nanosaniyeler mertebesinde gecikme farkları oluşur. $0111\to1000$
geçişinde, eğer örneğin en sağdaki bit diğerlerinden biraz daha yavaş değişirse, devre
kısa bir süreliğine $1001$ gibi hiç var olmaması gereken, yanlış bir ara değer üretebilir.

\begin{tanim}{Gray Kodu:}
Gray kodu, ardışık iki değer arasında \textbf{yalnızca tek bir bitin} değiştiği bir ikili
kodlama şemasıdır.
\end{tanim}

\textbf{Kodun adı nereden geliyor?} Gray kodu, adını Bell Telefon Laboratuvarları'nda çalışan
Amerikalı fizikçi ve mühendis Frank Gray'den (13 Eylül 1887-23 Mayıs 1969) alır. Purdue
Üniversitesi'nden 1911'de fizik diplomasıyla mezun olan Gray, kariyerinin büyük bölümünü
televizyon teknolojisine adadı: 1927'de erken televizyon sistemleri için "uçan nokta
tarayıcı" (flying spot scanner) fikrini önerdi, 1930'da ise iki yönlü, mekanik taramalı bir
televizyon sisteminin geliştirilmesine katkıda bulundu.

\begin{figure}[H]
\centering
\includegraphics[width=0.4\textwidth,keepaspectratio]{frank_gray.png}
\caption{Frank Gray, 1930'da geliştirdiği mekanik televizyon aygıtını incelerken.}
\label{fig:frank-gray}
\end{figure}

1947'de başvurduğu ve 1953'te kendisine verilen, "Pulse Code Communication" (Darbe Kodlu
Haberleşme) başlıklı 2,632,058 numaralı patentte Gray, analog bir sinyali sayısal darbelere
çeviren bir darbe kodu modülasyonu (PCM) aygıtı tanımladı. Bu aygıtın içindeki plaka,
ardışık değerler arasında yalnızca tek bir bitin değiştiği bir kod düzeniyle kaplanmıştı;
Gray, patent başvurusunda bu kodun o tarihte henüz tanınmış bir adı olmadığını belirtmişti.
Bugün bu kod, onun adıyla anılıyor.

\begin{figure}[H]
\centering
\includegraphics[width=0.85\textwidth,keepaspectratio]{frank_gray_patent.png}
\caption{Gray'in 2,632,058 numaralı patentinin ön sayfasından bir bölüm; yansıtılmış ikili
kodu taşıyan plakayı (15) içeren PCM tüpü (10) görülüyor.}
\label{fig:frank-gray-patent}
\end{figure}

Kodun kendisi aslında Gray'den çok önce de kullanılmıştı: Fransız mühendis Émile Baudot,
1870'lerde telgraf makinesinin harf çarkını sıralarken benzer bir kod düzeni kullanmış,
Alman-Avusturyalı Otto Schäffler de 1874'te benzer bir düzeni göstermişti. Ama kodun bugünkü
"Gray kodu" adı, onu titizlikle tanımlayıp somut bir mühendislik problemine (PCM, yani
darbe kodu modülasyonu ile yapılan analog-sayısal dönüşüm) uygulayan ve patentleyen
kişinin Gray olmasından geliyor.

Tablo~\ref{tab:gray-kodu}, $0$'dan $15$'e kadar düz ikili ile Gray kodunu karşılaştırmalı
gösterir. Her satırdan bir sonrakine geçerken Gray sütununda yalnızca tek bir bitin
değiştiğine dikkat edin.

\begin{table}[H]
\centering
\renewcommand{\arraystretch}{1.3}
\begin{tabular}{c@{\hspace{2em}}c@{\hspace{2em}}c}
\toprule
Ondalık & İkili & Gray Kodu \\
\midrule
0 & 0000 & 0000 \\
1 & 0001 & 0001 \\
2 & 0010 & 0011 \\
3 & 0011 & 0010 \\
4 & 0100 & 0110 \\
5 & 0101 & 0111 \\
6 & 0110 & 0101 \\
7 & 0111 & 0100 \\
8 & 1000 & 1100 \\
9 & 1001 & 1101 \\
10 & 1010 & 1111 \\
11 & 1011 & 1110 \\
12 & 1100 & 1010 \\
13 & 1101 & 1011 \\
14 & 1110 & 1001 \\
15 & 1111 & 1000 \\
\bottomrule
\end{tabular}
\caption{0'dan 15'e kadar ikili ve Gray kodu karşılaştırması.}
\label{tab:gray-kodu}
\end{table}

\begin{ornek}[$7\to8$ Geçişinde İkili ile Gray Kodunun Karşılaştırılması]
Düz ikilide $7\to8$ geçişi $0111\to1000$'dir: dört bitin dördü de değişir. Gray kodunda
aynı geçiş $0100\to1100$'dür: yalnızca en soldaki bit değişir, diğer üç bit aynı kalır.
Geçiş sırasında hangi an yakalanırsa yakalansın, okunan değer ya eski değer ya da yeni
değerdir. İki değerin karışımı olan geçersiz bir ara değer hiçbir zaman ortaya çıkmaz.
\end{ornek}

Gray kodunun tipik bir kullanım alanı, dönen bir milin açısal konumunu ölçen sensörlerdir
(encoder). Mil döndükçe, üzerindeki kodlanmış disk bir sensör tarafından okunur. Düz ikili
kullanılsaydı, birden fazla bitin aynı anda değiştiği geçiş anlarında yanlış okuma riski
olurdu; Gray kodu bu riski ortadan kaldırır.

\begin{bilginot}[İkiliden Gray Koduna Dönüşüm]
Yukarıdaki tabloyu ezbere bilmeye gerek yoktur; bir ikili sayıdan Gray koduna doğrudan bir
formülle geçilebilir. $n$ bitlik ikili sayı $b_{n-1}\ldots b_1b_0$ için, Gray kodunun
bitleri $g_{n-1}\ldots g_1g_0$ şöyle hesaplanır:
\[
g_{n-1}=b_{n-1}, \qquad g_i = b_{i+1}\oplus b_i \quad (0\le i\le n-2)
\]
Burada $\oplus$ sembolü \emph{XOR} (özel-veya) işlemidir; Bölüm~\ref{chap:bool-cebri}'de
tanıtılacaktır. Yani en soldaki bit aynen korunur, diğer her bit ise kendisinden bir
solundaki bitle XOR'lanır.
\end{bilginot}

\begin{ornek}[Formülle $7=(0111)_2$'nin Gray Koduna Çevrilmesi]
$b_3b_2b_1b_0=0111$ için:
\[
g_3=b_3=0,\qquad g_2=b_3\oplus b_2=0\oplus1=1
\]
\[
g_1=b_2\oplus b_1=1\oplus1=0,\qquad g_0=b_1\oplus b_0=1\oplus1=0
\]
Sonuç: $(0100)_2$. Tabloyla birebir uyuşuyor.
\end{ornek}

\begin{alistirma}[\kolay]
Yukarıdaki formülü kullanarak $(1010)_2$'nin ($10$'un ikili karşılığı) Gray kodunu
bulunuz.
\end{alistirma}

\noindent\textbf{Çözüm:} $b_3b_2b_1b_0=1010$ için:
\[
g_3=b_3=1,\qquad g_2=b_3\oplus b_2=1\oplus0=1
\]
\[
g_1=b_2\oplus b_1=0\oplus1=1,\qquad g_0=b_1\oplus b_0=1\oplus0=1
\]
Sonuç: $(1111)_2$. Tabloda $10$ satırına bakıldığında Gray kodu $1111$'dir. Doğrulandı.

\subsection{ASCII Karakter Kodu}
\label{sec:ascii}

Şu ana kadar gördüğümüz kodlar (BCD, Gray) hep \emph{rakamları} temsil ediyordu. Ama bir
bilgisayarın işlemesi gereken bilginin çoğu rakam değildir; harfler, noktalama işaretleri
ve klavyedeki her tuş da ikili düzende temsil edilmelidir.

Kaç bite ihtiyacımız var? Bölüm~\ref{sec:ikili-kodlar}'daki İkili Kod tanımını
hatırlayalım: $n$ bit, $2^n$ farklı eleman ayırt eder. Temsil etmemiz gereken kümeyi
sayalım: $26$ büyük harf, $26$ küçük harf, $10$ rakam, noktalama/özel karakterler
(\%, \$, @ gibi) ve birkaç görünmez kontrol karakteri (satır başı, sekme gibi). Bu toplam
$128$'e kadar çıkar. $2^6=64$ yetmez, $2^7=128$ yeterlidir; bu yüzden \textbf{ASCII $7$
bit} kullanır.

\begin{tanim}{ASCII (American Standard Code for Information Interchange):}
ASCII, $128$ karakteri (harfler, rakamlar, noktalama işaretleri ve kontrol karakterleri)
$7$-bit bir kodla temsil eden standart karakter kodudur. Bit dizisi, bu kitap boyunca
kullandığımız $0$-indeksli düzene uygun olarak $b_6b_5b_4b_3b_2b_1b_0$ şeklinde
adlandırılır; $b_6$ en soldaki (en anlamlı) bit, $b_0$ en sağdaki (en az anlamlı) bittir.
\end{tanim}

\begin{bilginot}[Tarihi Bir Not: ASCII'nin Orijinal Bit Numaralandırması]
ASCII standardının 1963 tarihli ilk hâli ve onu izleyen birçok ders kitabı, bitleri
$b_1$'den $b_7$'ye kadar ($b_0$ hiç kullanılmadan) adlandırır. Bu, telgraf/teleprinter
mühendisliğinden kalma eski bir numaralandırmadır. Modern dijital tasarımda (ve bu kitabın
geri kalanında, örn. Gray kodunda $b_{n-1}\ldots b_1b_0$) standart olan $0$-indeksli düzen
kullanılır; burada tutarlılık için bu düzeni tercih ediyoruz.
\end{bilginot}

Standart ASCII tablosu, $7$ biti ikiye bölerek okunur: en soldaki $3$ bit
($b_6b_5b_4$) tablonun \textbf{sütununu}, kalan $4$ bit ($b_3b_2b_1b_0$) tablonun
\textbf{satırını} belirler. Bu bölünme rastgele değildir: sütun $000$-$001$ tamamen
kontrol karakterlerine (görünmez, yazdırılamaz), sütun $010$-$011$ rakamlara ve temel
noktalamaya, sütun $100$-$101$ büyük harflere, sütun $110$-$111$ küçük harflere ayrılmıştır.
Yani sütun numarası, karakterin kategorisini anında belli eder.

\begin{table}[H]
\centering
\small
\begin{tabular}{@{}c|cccccccc@{}}
\toprule
$b_3b_2b_1b_0$ \textbackslash\ $b_6b_5b_4$ & 000 & 001 & 010 & 011 & 100 & 101 & 110 & 111 \\
\midrule
0000 & NUL & DLE & SP & 0 & @ & P & ` & p \\
0001 & SOH & DC1 & ! & 1 & A & Q & a & q \\
0010 & STX & DC2 & " & 2 & B & R & b & r \\
0011 & ETX & DC3 & \# & 3 & C & S & c & s \\
0100 & EOT & DC4 & \$ & 4 & D & T & d & t \\
0101 & ENQ & NAK & \% & 5 & E & U & e & u \\
0110 & ACK & SYN & \& & 6 & F & V & f & v \\
0111 & BEL & ETB & ' & 7 & G & W & g & w \\
1000 & BS & CAN & ( & 8 & H & X & h & x \\
1001 & HT & EM & ) & 9 & I & Y & i & y \\
1010 & LF & SUB & * & : & J & Z & j & z \\
1011 & VT & ESC & + & ; & K & [ & k & \{ \\
1100 & FF & FS & , & < & L & \textbackslash & l & \textbar \\
1101 & CR & GS & - & = & M & ] & m & \} \\
1110 & SO & RS & . & > & N & \textasciicircum & n & \textasciitilde \\
1111 & SI & US & / & ? & O & \_ & o & DEL \\
\bottomrule
\end{tabular}
\caption{Standart ASCII karakter tablosu (7 bit, $b_6\ldots b_0$).}
\label{tab:ascii}
\end{table}

\begin{ornek}[Harf 'A'nın ASCII Kodunun Bulunması]
Tabloda 'A' harfi, sütun $100$ ve satır $0001$'de yer alır. Yani $b_6b_5b_4=100$,
$b_3b_2b_1b_0=0001$, birleştirince: $(1000001)_2$.
\end{ornek}

Tabloya dikkatlice bakılırsa ilginç bir örüntü görülür: 'A' $=1000001$, 'a' $=1100001$.
Aralarındaki tek fark $b_5$ bitidir: $0$ ise büyük harf, $1$ ise küçük harf. Bu, klavye ve
yazılım tasarımını kolaylaştıran bilinçli bir tercihtir; büyük/küçük harf dönüşümü
donanımda tek bir bitin çevrilmesine indirgenir.

\begin{bilginot}[Neden 8 Bit (Byte) ile Saklanır?]
ASCII yalnızca $7$ bit gerektirdiği halde, bilgisayarlar veriyi $8$-bitlik gruplar
(\emph{byte}) halinde işler. Fazladan kalan $8$. bit boşa gitmez; genellikle
\textbf{parity} (hata sezme) için kullanılır. Bir sonraki alt bölümün konusu tam olarak
budur.
\end{bilginot}

\begin{alistirma}[\kolay]
Harf 'Z'nin ASCII kodunu, yukarıdaki tablodan sütun/satır konumunu bularak yazınız.
\end{alistirma}

\noindent\textbf{Çözüm:} 'Z' harfi, tabloda sütun $101$ ve satır $1010$'da yer alır. Yani
$b_6b_5b_4=101$, $b_3b_2b_1b_0=1010$, birleştirince: $(1011010)_2$.

\subsection{Hata Sezme: Parity Biti}
\label{sec:parity}

Bir önceki alt bölümün sonunda bıraktığımız yerden devam edelim: ASCII $7$ bit
($b_6\ldots b_0$) kullanır, ama bilgisayarlar veriyi $8$-bitlik byte'lar hâlinde işler.
Fazladan kalan $8$. bit genellikle \textbf{parity} için kullanılır. Bu bit, tam olarak
$b_7$'dir: $7$-bit ASCII ($b_6\ldots b_0$) artı $1$ parity biti ($b_7$), hiç boşluk
kalmadan tam bir byte oluşturur.

Bir veri, bir kablodan veya bir depolama ortamından geçerken, elektriksel gürültü
nedeniyle bir bit yanlışlıkla $0\to1$ veya $1\to0$ dönebilir. Parity biti, bu tür tek
bitlik hataları sezmek için kullanılan basit bir yöntemdir.

\begin{tanim}{Parity Biti:}
Parity biti, bir mesaja eklenen ve mesajdaki toplam $1$ sayısını istenen bir pariteye
(çift ya da tek) tamamlayan fazladan bir bittir.
\begin{itemize}[nosep, leftmargin=*]
    \item \textbf{Çift parity (even parity):} parity biti, toplam $1$ sayısı çift olacak
    şekilde seçilir.
    \item \textbf{Tek parity (odd parity):} parity biti, toplam $1$ sayısı tek olacak
    şekilde seçilir.
\end{itemize}
Sistemde ikisinden biri seçilir ve tutarlı biçimde kullanılır; çift parity daha
yaygındır.
\end{tanim}

\begin{ornek}[Harf 'A' İçin Parity Biti Hesaplanması]
'A' $=1000001$ (7 bit). İçindeki $1$ sayısını sayalım: $2$ tanedir (zaten çift).
\begin{itemize}[nosep, leftmargin=*]
    \item Çift parity isteniyorsa, parity biti $0$ kalır (toplam zaten çift): $01000001$.
    \item Tek parity isteniyorsa, toplamı tek yapmak için parity biti $1$ olur:
    $11000001$.
\end{itemize}
\end{ornek}

Alıcı tarafta kontrol basittir: alıcı, gelen $8$ bitteki $1$ sayısını sayar. Beklenen
pariteye (çift/tek) uymuyorsa, en az bir bitin yolda bozulduğu anlaşılır; bu durumda
yeniden gönderim istenir (NAK, \emph{Negative Acknowledge}, yani "olumsuz onay/red
bildirimi"). Uyuyorsa aktarım onaylanır (ACK, \emph{Acknowledge}, yani "onay
bildirimi").

\begin{bilginot}[Parity'nin Kritik Sınırlaması: Çift Sayıda Hatayı Yakalayamaz]
Eğer \textbf{tek sayıda} bit ($1$, $3$, $5$, $7$ tanesi) bozulursa, toplam $1$ sayısının
paritesi değişir ve hata sezilir. Ama \textbf{çift sayıda} bit ($2$, $4$, \ldots) aynı
anda bozulursa, toplam $1$ sayısındaki değişimler birbirini götürür; parite değişmez ve
hata sessizce kaçar.

Somut bir örnekle görelim: 'A'nın çift parity ile gönderilen hâli $01000001$'dir (içinde
$2$ tane $1$ vardır). Aktarım sırasında bu iki $1$ biti de $0$'a dönerse, alıcı
$00000000$ alır. Bu byte'taki $1$ sayısı $0$'dır, yine çift; parity kontrolü "hata yok"
der, oysa veri tamamen farklıdır.

Bunun temel nedeni, parity bitinin veri hakkında çok kaba bir özet tutmasıdır: yalnızca
toplam $1$ sayısının tek mi çift mi olduğunu bilir, hangi bitlerin değiştiğine dair hiçbir
bilgi taşımaz. İki bit birden değiştiğinde, birinin etkisi diğerininkini götürür ve
toplam parite üzerindeki net etki sıfırlanır.

Peki bu zayıflığa rağmen parity neden hâlâ yaygın kullanılır? Çünkü gerçek dünyadaki
gürültü kaynaklı hatalarda, aynı anda tam olarak iki (ya da dört, altı, \ldots) bitin
birden bozulması istatistiksel olarak çok daha nadirdir. En sık karşılaşılan hata türü
tek bir bitin bozulmasıdır, ve parity tam olarak bunu, tek bir ekstra bitle, ucuza yakalar.
Daha güçlü hata sezme/düzeltme gerektiren uygulamalarda (örn. bellek denetimi, ağ
protokolleri) checksum, CRC (döngüsel artıklık denetimi) veya Hamming kodu gibi daha
gelişmiş yöntemler kullanılır; bunlar bu kitabın kapsamı dışındadır.
\end{bilginot}

\begin{alistirma}[\kolay]
$A=0101100$ (7 bit) için çift parity bitini bulunuz.
\end{alistirma}

\noindent\textbf{Çözüm:} $0101100$ içindeki $1$ sayısı: $3$ tanedir (tek). Çift parity
istediğimiz için, toplamı çift yapacak parity biti gerekir: $3+1=4$ (çift), yani parity
biti $1$'dir.

Bu, Bölüm~\ref{sec:ikili-kodlar}'da başladığımız §1.7'nin son alt bölümüdür: BCD'den
Gray koduna, ASCII'den parity'ye kadar, sayıların ve sembollerin ikili düzende nasıl
kodlandığını ve bu kodlamalardaki hataların nasıl sezilebileceğini gördük.
